%%%%%%%%%%%%%%%%%%%%%%%%%%%%%%%%%%%%%%%%%%%%%%%%%%%%%%%%%%%
% letztes Commit:
% 359b2b49e8bd0487cda91e8035b629b36588231f
%%%%%%%%%%%%%%%%%%%%%%%%%%%%%%%%%%%%%%%%%%%%%%%%%%%%%%%%%%%
\ChapterCode{PSI}
\chapter
[\CO{[PSI]} Psychologie und psychische Betreuung]
{Psychologie und psychische Betreuung}
\label{chp:PSY}

\ChapterIntro
{%
~~~
}
{%
\begin{description}
  \item[Maintainer:] Sebastian Gabriel
  \item[Autoren:] Michael Auer, Sebastian Gabriel
  \item[Reviewer:] Michael Auer, Standard-Reviewprozess
  \item[Version:] {\DocVersion} ({\DocVersionInfo})
  \item[SHA1:] {\DocGitCommit}
\end{description}

{\TxTeilkapitelAASS}
}



\clearpage
\section{Anforderung, Beanspruchung/Belastung, Belastungsfaktoren}

\subsection{Anforderungen} 

Im Kapitel \E{Tätigkeitsbild} \CO{TAE} werden unter
\FullPrettyRef{topic:taetigkeitsbild-sanitaeter-belastung}
 die allgemeinen Anforderungen an Sanitätspersonal, wie
z.\,B. medizinisches Wissen und Kompetenz, Zuverlässigkeit, 
Eigenverantwortlichkeit, Fähigkeit Verantwortung
zu übernehmen, Einfühlungsvermögen, Teamfähigkeit und 
\EE{Frustrationstoleranz} besprochen. Im Folgenden wird auf
die damit einhergehende Belastung durch diese und andere
Anforderungen und die daraus resultierenden Belastungen
eingegangen.
% \A{fachliches Wissen, praktische Kompetenzen, dieses in
% schwierigen Situationen umsetzen zu können professionelle 
% Haltung, social/soft skills zwischenmenschliche Fähigkeiten,
% Frustrationstoleranz,  Umgang mit mehrdeutigen Situationen, 
% \enquote{schwierigem Klientel}, Teamdifferenzen.}
% 
% Die Tätigkeit im Rettungs- und Krankentransportwesen geht
% aufgrund der Komplexität der Aufgaben mit einer Fülle von
% körperlichen und psychischen Anforderungen einher. Einen nur
% kargen Hinweis darüber, welche Qualitäten von einem
% (Zivildienstleistenden) Sanitäter erwartet werden, liefert
% das Anforderungsprofil des AMS, nämlich:
% \begin{itemize}
%   \item Körperliche und psychische Belastbarkeit +
%   \item rasches Auffassungs- und Reaktionsvermögen, +
%   item  Einfühlungsvermögen +
%   \item  Beobachtungsgabe, +
%   \item  Kommunikations- und Teamfähigkeit, +
%   \item Zuverlässigkeit, +
%   \item Verantwortungsbewusstsein +
% \end{itemize}
% Im Folgenden soll der Versuch unternommen werden diese
% nüchternen Auflistung von Begrifflichkeiten plastischer
% darzustellen und zu erweitern.
% 
% Es versteht sich von selbst, dass die professionelle
% Versorgung erkrankter bzw. verunfallter Menschen
% medizinisches Grundwissen voraussetzt und damit auch
% einhergehend die praktischen Kompetenzen, dieses
% \EE{zuverlässig und eigenverantwortlich} in schwierigen
% Situationen umzusetzen zu können. 
% 
% Dringliche
% Versorgungssituationen bzw. Notfälle, in denen jede
% zeitliche Verzögerung zulasten des Patienten fallen könnte,
% machen ein \EE{schnelles Einschätzen der Lage} so wie die
% Bereitschaft \EE{Verantwortung zu übernehmen} notwendig. Da
% es sich bei Versorgungssituationen stets um eine Interaktion
% zwischen zwei Personen -- nämlich Versorgendem und zu
% Versorgenden -- handelt, gelten zwischenmenschliche
% Fähigkeiten wie höfliche Umgangsformen,
% \EE{Kommunikationsfähigkeit} sowie \EE{Einfühlungsvermögen}
% als Grundvoraussetzung. 
% 
% Erkrankung bzw. Unfall stellen für
% die Betroffenen in häufigen Fällen kein alltägliches
% Geschehen dar, sondern haben oft einen
% kritisch-krisenhaften, jedenfalls aber
% frustrierend-belastenden Charakter, weshalb an dieser Stelle
% die besondere Wichtigkeit einer \EE{einfühlenden und
% verständnisvollen Grundhaltung} betont werden soll. Dies
% kann für den professionellen Helfer auch damit einhergehen, seine
% eigenen Bedürfnisse bzw. Empfindungen für den Moment
% hintanzustellen, also eine hohe \EE{Frustrationstoleranz}
% erforderlich machen.  
% 
% Im Rettungsdienst beschäftigte
% Personen arbeiten grundsätzlich die meiste Zeit in Teams,
% sodass persönliche Qualitäten wie Toleranz, Kritikfähigkeit
% und Selbstdisziplin (\EE{\enquote{Teamfähigkeit}}) ebenso von hoher
% Bedeutung sind.
% 
% Weiters sei auf die Wichtigkeit \EE{physischer
% Belastungsfähigkeit} als erforderliches Kriterium
% hingewiesen: das Tätigkeitsfeld des Rettungssanitäters
% umfasst diverse Aufgaben, die mit körperlichem Einsatz und
% dementsprechender Anstrengung verbunden sind. 
% 
% Zu guter Letzt
% und im besonderen sei auf die Bedeutung \EE{psychischer
% Belastbarkeit} hingewiesen, die auch in den folgenden
% Kapiteln zum Thema gemacht werden soll, v.a. in Hinblick auf
% den Umgang mit psychischem Stress: 
% 
% Sanitätsdienst zu
% leisten, bedeutet die nahezu tägliche Konfrontation mit dem
% persönlichen Leid von Mitmenschen, des weiteren die
% Verantwortlichkeit – ggf. auch noch unter Zeitdruck – darauf
% professionell zu reagieren und zusätzlich die besonderen und
% generellen Stressauslöser des Arbeitssettings
% \enquote{Sanitätsdienst} zu meistern. Dies bedingt eine hinreichend
% funktionierende psychische Verarbeitungsfähigkeit, um ein
% gesundes seelisches Gleichgewicht wahren und pathologischen
% Entwicklungen (Burnout, Belastungsreaktion etc.) präventiv
% entgegenwirken zu können.

\subsection{Beanspruchung/Belastung}
\A{(Kurze Einführungsphase)   Parallelen
(Beanspruchungsdauer, -intensität) bzw. Unterschied zw.
Mechanischer und psychischer Beanspruchung = interpersonelle
Faktoren (Wahrnehmung, Coping, Expertise) vs. strukturelle
Gegebenheiten (s.u.). }
\label{topic:beanspruchung-psi}
\label{topic:belastung-psi}

Unter dem Begriff \T{Beanspruchung} versteht man in der
Technik \Speech{\enquote{die Auswirkung einer äußeren Belastung auf
das Innere eines Körpers bzw. eines Werkstoffes}}
\cite{wiki:beanspruchung20110814}. Mechanische Beanspruchung
ist also durch die Ausbildung von mechanischen Spannungen
gekennzeichnet. Analog dazu kann man psychische
Beanspruchung als die \EE{Entstehung psychischer
Spannungszustände} verstehen, die im größeren Ausmaß als
psychische Belastung bzw.
Stresserleben bezeichnet werden. 


Reaktionen auf eine
Belastung lassen sich in verschiedene
Ebenen aufteilen: 
\begin{itemize}
  \item die körperliche \Z{(physiologische)},
  \item die gedankliche \Z{(kognitive)},
  \item die gefühlsmäßige \Z{(emotionale)} sowie die
  \item Verhaltensebene \Z{(behaviorale)}. 
\end{itemize}



Umstände wie
\EE{Beanspruchungsdauer} (kurz -- lang) bzw. 
\EE{-intensität} (stark -- schwach) haben einen wesentlichen
Einfluss auf das Beanspruchungsniveau. In der Psyche
spielen jedoch -- im deutlichen Gegensatz zur Technik --
mehrere \EE{vermittelnde Faktoren} eine wesentliche Rolle: 

\begin{itemize}
  \item Persönliche Wahrnehmung,  
  \item kognitive Bewertungen, 
  \item Verarbeitungsstrategien und Abwehrmechanismen sowie
  \item Routine, Vorerfahrung und Wissen. 
\end{itemize}



Diese können in
Beanspruchungssituationen bewusst oder unbewusst zum Tragen
kommen und sollen im Folgenden besprochen werden.
% Inhaltliche und strukturelle Aspekte des Sanitätsberufs
% stellen diesbzgl. wesentliche Einflussfaktoren dar und
% sollen im Folgenden erläutert werden.

\subsection{Belastungsfaktoren}
\A{- Mehrdeutige Situationen erhöhen erlebte Beanspruchung 
- Ruhe- und Alarmphasen - workflow, 
- Frustrane Erlebnisse (Vorwürfe von Patienten/Angehörige), 
- Erlebtes arbeitet emotional und kognitiv weiter,
- Rollenkonflikte, 
- persönlich adäquate Balance zwischen professioneller verständnisvoller Haltung, strukturellen bzw. fachlichen Notwendigkeiten und persönlichen Grenzen u. Bedürfnissen.}

 

Belastungsfaktoren kann man in die
Bereiche Aufgabenstruktur, Rollen- und Interaktionsstruktur
einteilen.\ZFootnote{\Study{Bengel} \cite{BengelPsyNotfallmedRd2-PsyBelRdPers}}


\TextSlide{Aufgaben}
{\begin{itemize}
  \item \EE{Arbeitsumgebung}: Sanitäter sehen sich häufig
  mit \EE{komplexen, mehrdeutigen Situationen} konfrontiert, wie
  z.\,B. bei Notfällen bzw. Erkrankungen, die die
  Handlungsplanung bei der Versorgung erschweren und damit
  nicht zuletzt auch die emotionale Anspannung erhöhen können.
  Auch \EE{Faktoren der Arbeitsumgebung} haben darauf einen
  Einfluss: der Einsatzort ist meist unbekannt, widrige
  Umstände wie erschwerte Zugänglichkeit, behindernde
  Schaulustige, emotionalisierte Angehörige, Dunkelheit oder
  schlechte Wetterbedingungen spielen eine Rolle.
  
  \item \EE{Workflow}: Während des Diesntes wechseln
   Ruhe- und Bereitschaftsphasen mit plötzlichen Aktivitäts-
   und Alarmphasen. Langfristig kann dieser \EE{unvorhersehbare bzw. unkontrollierbare
  Wechsel im Arbeitsrhythmus} als unangenehm wahrgenommen
  werden und zu Stresserleben beitragen.
  
  \item \EE{Psychologische Anforderung}: Nicht zuletzt
  stehen Mitarbeiter helfender bzw. versorgender Systeme nicht
  nur körperlich sondern auch psychisch verletzten Menschen
  gegenüber, was zur Folge haben kann, dass sie sich mit
  \EE{leidenden Äußerungen}, aber unter Umständen auch mit
  \EE{verbalen Vorwürfen und der Kritik Betroffener bzw.
  Angehöriger} konfrontiert werden.
  
  \item \EE{Nachwirkung}: Eine Besonderheit psychischer
  Beanspruchung stellt die Tatsache dar, dass \EE{emotionale
  bzw. gedankliche Verarbeitung} stressreicher Situationen
  bewusst oder unbewusst \EE{nachwirkt}. Wiederaufkeimende
  Gedanken und Gefühle an vergangene Einsätze halten sich
  nicht an Dienstzeiten, sodass auch nach Dienstschluss
  innerlich emotional ge- und verarbeitet wird.
\end{itemize}}
{\begin{itemize}
  \item {Arbeitsumgebung}
  \item {Workflow}
  \item {Psychologische Anforderung} 
  \item {Nachwirkung}
\end{itemize}}
{}
{}
{}

\Text{Rolle}
{\begin{itemize}
  \item \EE{Rollenkonflikt}: Forderungen/Erwartungen können
  nicht erfüllt werden, da der Arbeitgeber/Vorgesetzte die
  Kompetenzen des Personals beruhend auf gesetzlichen
  Bestimmungen bzw. ökonomischen Ressourcen festlegt,
  gleichzeitig das Personal seinen Kompetenzbereich aber
  aufgrund eigener Interpretationen anders auslegt und
  z.\,B. in einer Notfallsituation entsprechend eigenem
  Rollenverständnis handelt.
  
  \item \EE{Rollenambiguität} (Unsicherheit/Mehrdeutigkeit):
  Sie entsteht, wenn Informationen zur Aufgabendurchführung nur
  unzureichend zur Verfügung stehen, z.\,B. wenn
  verschiedene Beteiligte (Arbeitgeber, Ausbilder und
  Gesetzgeber) zu einer lebensrettenden Maßnahme
  unterschiedliche oder keine verbindlichen Aussagen treffen.
  
  \item \EE{Rollenüberforderung}: Mitunter wird 
  mehr von einer Person gefordert, als sie in der ihr
  verfügbaren Zeit oder mit den ihr zu Verfügung stehenden
  Mitteln erreichen kann, die Erwartungshaltung \E{muss}
  zwangsläufig enttäsucht werden. Dies kann z.\,B.
  durch die Patienten selbst, aber auch durch Dritte wie Angehörige
  geschehen. Ebenso muss das
  ersteintreffende Behandlungsteam bei
  Großschadensereignissen seine Handlungsschritte/Ressourcen
  nach Prioritäten ordnen und kann nur begrenzt Hilfe
  leisten.
\end{itemize}}
{\begin{itemize}
  \item {Rollenkonflikt}
  \item {Rollenambiguität}
  \item {Rollenüberforderung}
\end{itemize}}
{}
{}
{}

\Text{Rolle}
{
\begin{itemize}
  \item
  Sanitätsdienst als berufliche Rolle geht mit diversen
  Anforderungen seitens der eigenen Person, der Organisation
  und nicht zuletzt der Gesellschaft einher. Demgegenüber ist
  er durch bestimmte Rahmenbedingungen wie zeitliche,
  personelle, strukturelle Ressourcen und gesetzliche
  Bestimmungen beschränkt. Daraus können Situationen
  entstehen, die zu
  \item \EE{Rollenkonflikt} (z.B. Konflikt
  zwischen dem eigenen Anspruch und der Ressourcenlage der
  Organistation),
  \item \EE{Rollenambiguität}  (z.B. unklare
  Ausgangslage bzgl. lebensrettender Maßnahme aufgrund
  widersprüchlicher Lehr- bzw. Gesetzesmeinung) und auch
  \item
  \EE{Rollenüberforderung} (z.B. Überschreitung der
  persönlichen oder zeitlichen Ressourcen bei
  Großschadensereignis oder aufgrund zu hoher
  Erwartungshaltung seitens des Betroffenen bzw. deren
  Angehörige)
\end{itemize}
}
{\begin{itemize}
  \item {Rollenkonflikt}
  \item {Rollenambiguität}
  \item {Rollenüberforderung}
\end{itemize}}
{}
{}
{}



Zu guter Letzt sei daraufhingewiesen, dass die individuell
erlebte Beanspruchung auch von der Fähigkeit des Einzelnen
abhängt, in seiner Rolle als in einem Versorgungssystem 
tätige Person \EE{eine persönlich angemessene Balance}
zwischen einer \EE{professionellen verständnisvollen Haltung
(z.\,B. soft skills wie empathisches Umgehen etc.),
strukturellen bzw. fachlichen Erfordernissen seiner
Rollenfunktion als Sanitäter (z.\,B. Dienstvorschriften,
Behandlungsstandards) und persönlichen Grenzen u.
Bedürfnissen (z.\,B. eigenes Wohlbefinden, Gefühle) zu
finden (Beispiel siehe Exkurs zu Rollenambivalenz)}.

Die individuelle erlebte Beanspruchung ist daher auch von der Fähigkeit des Einzelnen abhängig in dieser Rolle eine persönlich angemessene Balance zwischen den eignen Ansprüchen bzw. Ansprüchen von außen(verständnisvolle Grundhaltung, Dienstvorschriften, Behandlungsstandards) und persönlichen Grenzen und Bedürfnissen (persönliches Wohlbefinden, Gefühlserleben) zu finden.

\Text{Exkurs: Rollenambivalenz}
{\begin{center}
\TableBaseModification
\begin{tabularx}{\linewidth}{XcX}
\toprule\addlinespace
Professionist 
&
vs.
& 
Person 
\\\addlinespace
Verantwortlichkeit 
& 
vs.
& 
pers. Gefühle
\\\addlinespace
Schnelles, aber professionelles und möglichst fehlerfreies
Handeln bei gleichzeitig hoher Verantwortung
&
vs.
&
Emotional berührende Situationen, Konfrontation mit den
Themen Unfall/Krankheit, Leid, Trauer, Schock und Tod
\\\addlinespace\bottomrule
\end{tabularx}
\end{center}


% Sanitäter sehen sich in ihrer Rollenfunktion mit einem
% breiten Spektrum an belastenden Themen konfrontiert, deren
% erschöpfende Darstellung hier nicht abgehandelt werden kann.
% Exemplarisch soll der Rollenkonflikt, der sich aus der
% Auseinandersetzung mit Belastungen im beruflichen Kontext
% ergibt, erörtert werden.

\subparagraph{Exemplarisch: Verantwortlichkeit und
persönliche Gefühle}

Die Rolle als Sanitäter beinhaltet, dem
\E{Versorgungsauftrag} (Versorgung erkrankter Menschen) 
verpflichtet zu sein, d.\,h. sich für die Gesundheit bzw. in letzter Instanz evtl. für
das Leben eines anderen Menschen verantwortlich zu fühlen.

\Z{Im Rahmen dieser Aufgabe werden dem professionellen
Helfer diverse Handlungs- und Entscheidungskompetenzen
überantwortet -- z.\,B. die Wahl der Maßnahmen am Patienten,
die Organisation des Einsatzes etc. --, zwischen denen er
womöglich auch unter Zeitdruck nach bestem Wissen und
Gewissen entscheiden soll.

Gleichzeitig übt eine fühlende Person diese
verantwortungsvolle Versorgungsrolle aus, die sich mit
berührenden Themen wie Unfall/Krankheit, Leid, Trauer,
Schock und Tod konfrontiert sieht, sodass ihr Verhalten und
Erleben beeinflusst wird durch Gefühle der Anteilnahme,
Belastungs- und Frustrationsgefühle, persönliche
Erinnerungen, Stimmungslagen etc.}

Diesen Rollenkonflikt in einer angenehmen und funktionellen
Balance zu halten setzt die Bereitschaft zur Selbstreflexion
und Auseinandersetzung mit persönlichen Grenzen und
Emotionen voraus. Diesbzgl. Ungleichgewichte  können sich
einerseits durch das übermäßiges Verhaften in der
Professionistenrolle (\Speech{\enquote{Versorgungsroboter}}),
andererseits aber auch durch \E{zu hohes Einbringen}
persönlicher Gefühlsanteile aufgrund von Gefühlsansteckung (=Übernahme
der Gefühle des Betroffenen) bis hin zur Lähmung  durch
emotionalen Kontrollverlust (z.\,B. emotionaler Schock)
ausdrücken.

Es muss betont werden, dass eine ausgeglichene
Rollenaufteilung \enquote{Sanitäters vs. Person} einen
Schutzfaktor (Ressource) darstellt, da sie hilft das persönliche vom
professionellen zu trennen, aber gleichzeitig auch erlaubt
seinen persönlichen Stil in die Versorgung
einzubringen. Dies erfordert jedoch eine ausgewogene
Integration der beiden dargestellten Aspekte
\enquote{Sanitäterrolle} und \enquote{Rolle als emotionales Individuum}.
Das Bewusstmachen persönlicher Gefühle und Grenzen -- z.\,B.
durch \EE{Selbstreflexion}, Gespräche etc. -- ist hierbei
wesentlich und hilfreich.}
{}
{}
{}
{}



\Text{Interaktion}
{Sie betrifft die \EE{Kommunikation und Kooperation zwischen
Vorgesetzten und Mitarbeitern sowie zwischen Kollegen}
untereinander: \EE{negative zwischenmenschliche Beziehungen}
zwischen Kollegen oder Vorgesetzten – gekennzeichnet durch
mangelndes Vertrauen, geringe Unterstützungsbereitschaft und
geringes Interesse zuzuhören bzw. sich mit den Problemen
anderer auseinanderzusetzen –, \E{Hierarchieprobleme} sowie
\E{mangelnde Rückmeldung über Einsatzerfolge} zählen laut
einer Studien\ZFootnote{Gigsby \& Mc New, 1988 zitiert nach
\cite{BengelPsyNotfallmedRd2}} zu den
Hauptbelastungsfaktoren, ebenso wie beeinträchtigte
Kooperationsmöglichkeiten und \E{unzureichende Anerkennung}
sowohl durch Mitarbeiter als auch Leiter \ZFootnote{Udris,
1981 zitiert nach \cite{BengelPsyNotfallmedRd2}}. Demzufolge
lässt sich dem \EE{sozialen Zusammenhalt} im Rettungsdienst
für die
erfolgreiche Bewältigung der Aufgaben eine hohe Bedeutung
zuschreiben.}
{
\begin{itemize}
  \item Kommunikation und Kooperation zwischen
  Vorgesetzten und Mitarbeitern sowie Kollegen
  \item Sozialer Zusammenhalt
  \item
\end{itemize}
}
{}
{}
{}








\section{Stress: Ursachen und Einflussfaktoren der Stressentstehung}

\label{topic:stress-ursachen}

\subsection{Stress -- Ein Defintionsversuch?}

\label{topic:stress}

\Definition{Körperliche und psychische
Reaktionen, die durch die
Wahrnehmung bestimmter Reize in unserer Umwelt bzw.
innerhalb einer Person ausgelöst werden können \cite{ZimbardoPsych6}.}

% In der Alltagssprache mag der Begriff \enquote{Stress} als ein
% Phänomen mit einer sehr negativen Bedeutung Verwendung
% finden. Wissenschaftlich betrachtet können unter diesem
% Begriff allgemein \EE{physiologische und psychische
% Reaktionen} zusammengefasst werden, \EE{die durch die
% Wahrnehmung bestimmter Reize in unserer Umwelt} \A{(siehe
% auch SOC-Modell)} \EE{bzw. innerhalb einer Person ausgelöst
% werden können} (=Stressoren; Zimbardo, 1995). 

Das Phänomen  Stress kann als Summe der
Reaktionen auf die bewusste und/oder unbewusste Wahrnehmung
äußerer Umweltreize oder intrapersoneller Reize verstanden
werden und sich auf vielen Ebenen (kognitiv-emotionaler und
Verhaltensebene) in unterschiedlicher Weise manifestieren.
Aufgrund seiner Funktion, den Organismus auf eine Handlung
vorzubereiten/in \EE{Alarmbereitschaft} zu versetzen, ist er
per se nicht als positiv oder negativ zu bewerten, sondern
stellt eine \E{überlebenswichtige, evolutionär entstandene
natürliche Reaktion} dar.


\ZFootnote{In jeder
Sekunde strömen ca. 11 Millionen Reize auf unser
Nervensystem ein. Sie werden zunächst einer unbewussten
Sofortanalyse unterzogen. Dabei findet v.\,a. eine
kognitiv-emotionalen Bewertung statt. 

V.\,a. neue, mit
vorhandenen Mustern nicht übereinstimmende Reize werden
bewusst und sorgfaltig verarbeitet, indem wir ihnen
Aufmerksamkeit widmen und sie zur weiteren Bearbeitung im
Kurzzeitgedächtnis zwischengespeichert werden. Sowohl die
\EE{bewusst}, als auch die \EE{unbewusst} (subliminal,
unterschwellig, also tachistoskopisch, d.\,h. \VonBis{1}{5}\MS* lang
dargebotene Wahrnehmungen, Priming) \EE{verarbeiteten Reize}
wirken sich auf das Verhalten aus (Verhaltenswirksamkeit).
Dabei herrscht das Primat des Affekts, d.\,h. Gefühle
manifestieren sich oftmals vor der bewussten Wahrnehmung,
weil die Verarbeitung in limbischen Arealen (Amygdala, Teil
des Stammhirns = mit unbewussten Prozessen assoziiert) der
kortikalen Verarbeitung (Großhirnrinde = mit bewusster
Wahrnehmung assoziiert) vorausgeht. Dieser evolutionär
entstandene Mechanismus \EE{bereitet} als
\EE{Frühwarnsystem} den Organismus \EE{auf eventuellen Kampf
oder Flucht} (fight or flight) vor. Oberflächlich betrachtet
stellt Stress also eine \EE{Aktivierung des Organismus} dar,
die ihm helfen soll ausreichende \EE{Mittel zur Bewältigung
einer Situation/eines Problems zur Verfügung zu stellen} –
was per se also nicht als gut oder schlecht bewertet werden
kann. }

\begin{figure}[H]
    \begin{center}
    \includegraphics[width=\linewidth]{./images/hirtler-aass/Parasymp-Symp-Cartoon.pdf}
    \Captionof{figure}{Stress hat, genauso wie Ruhephasen,
    eine wichtige Funktion. \SourcePicture{Hirtler}}
    \end{center}
\end{figure}

\Commentary[Stress]{Da es von Stress so viele Definitionen
geben mag, wie Forscher die sich dem Phänomen gewidmet haben, sei hier auf
eine erschöpfende Auflistung solcher verzichtet und
stattdessen exemplarisch die wichtigsten Stressmodelle
dargestellt. Es soll betont werden, dass die vorgestellten
Theorien bewusst wenig detailreich und in
nicht-chronologischer Reihenfolge ihrer Entstehung
dargestellt sind, um das Verständnis zu erleichtern.}





\subsection{Stressmodelle}

\subsubsection{Modell der Homöostase-Allostase}


Unter \T{Homöostase} versteht man ein (fiktives
biopsychosoziales) \E{Gleichgewicht eines Individuums}; Stress
kann dieses Gleichgewicht gefährden.\ZFootnote{Unter 
\T{Allostase} versteht man einen \E{Regelkreis,
dessen Sollwert stärker verstellbar ist} als bei
homöostatischen Systemen; trotz äußerer Störfaktoren wird
das innere Milieu dank verschiedener
\E{Anpassungsreaktionen} konstant gehalten; eine
Fehlregulation allostatischer Systeme durch zu \E{häufigen,
intensiven oder lange andauernden Stress (z.\,B. chronischer
Stress) kann pathogen wirken}, aufgrund langfristiger
physiologischer Veränderungen.
% (s. Grafik unten) 
}



Mögliche
Belastungsquellen, die die Anpassungsmöglichkeiten des
Organismus  überfordern und krankheitsauslösend wirken
können, wenn sie als unkontrollierbar erlebt werden sind
z.\,B.
\begin{itemize}
  \item \EE{kritische Lebensereignisse} (\EE{life events})
  z.\,B. der Verlust nahe stehender Menschen, Scheidung,
  Unfall, schwere Erkrankung, Verlust des Arbeitsplatzes,
  Schwangerschaft etc.
  \item \EE{außergewöhnlicher/traumatisiserender Stress}
  z.\,B. Katastrophen, lebensbedrohliche Situationen,
  Verbrechensopfer, Gewalterfahrungen, emotionaler
  Missbrauch, Folter, Gefangenschaft.
  \item \EE{alltägliche Hintergrundstressoren} (\EE{daily
  hassles}), die sich mit der Zeit aufsummieren persönliche
  Stressoren z.\,B. dysfunktionale Gedanken wie
  Selbstentwertung, überhöhte Anforderungen an sich selbst
  \item \EE{Idiosynkratische Körperreaktionen} auf Stress
  \item \EE{defekte Feedbackschleifen}, die die Homöostase
  aus dem Gleichgewicht bringen.
\end{itemize}


\Fazit{Postuliert wird ein 
Gleichgewicht, das durch Stressoren beeinflusst wird.
Anpassungsreaktionen versuchen ein
dadurch entstandenes Ungleichgewicht auszugleichen.
Kritische Lebensereignisse, chronische oder besonders
intensive Stressoren können zu Fehlregulationen und in
weiterer Konsequenz zu pathologischen Veränderungen führen.}

\subsubsection{Gleichgewichtsmodell}

Das Gleichgewichtsmodell geht von einem angenommenen biopsychosozialen Gleichgewicht aus nach dem jedes Individuum strebt, d.h. sowohl organisches als auch psychische und zwischenmenschliche Prozesse sind darauf ausgerichtet ein als angenehm erlebtes Gleichgewicht zu erreichen und aufrechtzuerhalten (=Anpassungsprozesse). Stress kann das persönliche Gleichgewicht gefährden und Anpassungsreaktionen erforderlich machen.

Folgende Belastungsquellen können die Anpassungsmöglichkeiten des
Organismus  überfordern, beeinträchtigen und krankheitsauslösend wirken, v.a. wenn sie als unkontrollierbar erlebt werden:

\begin{itemize}
  \item \EE{kritische Lebensereignisse} (\EE{life events})
  z.\,B. der Verlust nahe stehender Menschen, Scheidung,
  Unfall, schwere Erkrankung, Verlust des Arbeitsplatzes,
  Schwangerschaft etc.
  \item \EE{außergewöhnlicher/traumatisiserender Stress}
  z.\,B. Katastrophen, lebensbedrohliche Situationen,
  Verbrechensopfer, Gewalterfahrungen, emotionaler
  Missbrauch, Folter, Gefangenschaft.
  \item \EE{alltägliche Hintergrundstressoren} (\EE{daily
  hassles}), die sich mit der Zeit aufsummieren persönliche
  Stressoren z.\,B. dysfunktionale Gedanken wie
  Selbstentwertung, überhöhte Anforderungen an sich selbst
  \item \EE{Idiosynkratische Körperreaktionen} auf Stress
  \item \EE{defekte Feedbackschleifen}, die die Homöostase
  aus dem Gleichgewicht bringen.
\end{itemize}

Das Diathese-Stressmodell geht davon aus, dass Menschen über eine individuelle tw. angeborene, tw. im Laufe ihrer Lebenserfahrung erworbene Grundausstattung (Repertoire an psychischen oder im Verhalten liegenden Anpassungsmöglichkeiten) verfügt, Stress zu begegnen und somit anfälliger oder weniger anfällig für Stressfolgen sind (→ individuelle Anfälligkeit bzw. Verletzlichkeit für Stress = Diathese). 

Im Modell des Allgemeinen Adaptationssyndroms werden Stressreaktionen in drei Phasen unterteilt: eine Alarmphase, in der körperliche Prozesse einsetzen um Handlungsbereitschaft zu erzeugen(Erhöhung von Blutzucker, Herz- und Atemfrequenz). Es folgt die Widerstandsphase, in der es zur Aktivierung aller Energiereserven kommt, um Handlungen setzen zu können. Werden sämtliche Energiereserven aufgebraucht, setzt anschließend die Erschöpfungsphase ein (Überforderung durch chronischen Stress, Schwächung der Immunabwehr) – eine Erholungsphase ist erforderlich um Energie-Ressourcen "aufzufüllen" und neuerlichen Wiederstand zu ermöglichen. Insgesamt wird eine Reaktions-/Alarmbereitschaft erzeugt, die ein adäquates Reagieren auf den Stressor ermöglichen soll, solange
Ressourcen vorhanden sind. Die verschiedenen Phasen sind mit jeweiligen
physiologisch-vegetativen Reaktionen verknüpft. 

Das neuroendokrine Modell ergänzt das bisher beschriebene dahingehend, dass verschiedene Stressoren individuell je nach Situation emotional unterschiedlich bewertet werden können (Furcht, Ärger, Niedergeschlagenheit etc.) und damit spezifische biochemische Reaktionen folgen. D.h. je nach emotionaler Bewertung folgen entsprechende biochemische Prozesse.

Auch im dzt. geläufigen Transaktionalen Stressmodell spielen Bewertungsprozesse eine wesentliche Rolle: es unterscheidet eine primäre und sekundäre Bewertung. Im ersten Bewertungsschritt wird geprüft „Ist ein Reiz oder eine  Situation für mich überhaupt relevant?“. Falls ja, folgt die Frage „Stellt er/sie eine Herausforderung oder Bedrohung für mich da?“ – hier spielen die Einschätzung des eigenen Könnens und der Aktivationsgrad eine wichtige Rolle. In der zweiten Bewertungsphase wägt der Betroffene seine Handlungsmöglichkeiten ab. Erlebt sich der/die Betroffene als handlungsfähig kann er /sie entsprechend reagieren und zur Verarbeitung der Situation (Coping)mit dem Ziel der Bewältigung (durch Verhalten, gedanklich oder emotional) übergehen. Im Anschluss erfolgt eine Neubewertung der Situation. Die Einschätzung keine Handlungsmöglichkeit zu haben erzeugt Stresserleben. 
 
\subsubsection{Allgemeines Adaptationssyndrom}

Nach initialer Alarmphase und passagerer
Widerstandsphase mit Mobilisierung der Reserven tritt die
Erschöpfungsphase ein, weil sich die Energie-Reserven
verbrauchen. Die verschiedenen Phasen sind mit jeweiligen
physiologisch-vegetativen Reaktion verknüpft. Insgesamt wird
eine Reaktions-/Alarmbereitschaft erzeugt, die ein adäquates
Reagieren auf den Stressor ermöglichen soll, solange
Ressourcen vorhanden sind. 
\Commentary[Allgemeines
Adaptationssyndrom]{Das Allgemeine
Adaptationssyndrom (Von Selye postuliert, zitiert nach
\cite{ZimbardoPsych6}) ist als eine komplexe Antwort des
Körpers zu verstehen, wenn die Homöostase des Organismus
durch Stressoren wie z.\,B. Gefahr, Lebensbedrohung,
Belästigung etc. gestört wird. Es gliedert sich in drei
Phasen:


\EE{Alarmphase}: diese ist unterteilt in Schock- und
  Gegenschockphase. Schockphase: Herzfrequenz erhöht sich, Blutdruck und Blutglucosespiegel sinken.
  Gegenschockphase: Sympathikus wird aktiviert und ACTH ausgeschüttet.

\EE{Widerstandsphase}: Aktivierung aller
  Energiereserven; erhöhte Aktivierung des Sympathikus,
  Cortisol-Ausschüttung nimmt zu (wg. ACTH-Freisetzung in
  der Alarmphase; kann zu einer kompensatorischen
  Hypertrophie der Nebennierenrinde führen), Veränderung
  vegetativer Parameter;  Daueraktivierung des Sympathikus
  kann das Risiko von Herz-Kreislauf-Erkrankungen erhöhen.

\EE{Erschöpfungsphase}: Aufbrauchen der Reserven,
  Dekompensation der Stressreaktion bei chronischem Stress,
  Immunsuppression aufgrund dauernder Cortisol-Ausschüttung,
  Wachstum/Reparatur und Reproduktionsfunktion sind gehemmt.
  Eine Erholungsphase ist erforderlich um Ressourcen
  \enquote{aufzufüllen} und Wiederstand zu ermöglichen.

}

\A{ZIMBARDO!!!}


\subsubsection{Psychoneuroendokrines Stressmodell}

Je nach Stressor und
dessen emotionaler Bewertung (Furcht, Ärger, Depression) initiiert
der Körper spezifische biochemische
Reaktionen\ZFootnote{Psychoneuroendokrines Stressmodell: Von
Henry (1986) postuliert, stellt es eine Differenzierung des
Selye-Modells dar}. 
\Commentary[Psychoneuroendokrines Stressmodell]{Zuerst
erfolgt eine Verarbeitung im frontotemporalen Kortex
(Aufmerksamkeit), dann erfolgt die emotionale Bewertung:
\EE{Furcht}: Anstieg des Adrenalin \Z{(basale Amygdala)},
\EE{Ärger}: Anstieg des
  Noradrenalin und Testosteron \Z{(zentrale Amygdala)},
\EE{Depression}: Anstieg des Cortisol,
  Abnahme des Testosteron \Z{(Hippocampus)}
}

\Fazit{Die Psychophysiologische Reaktion  auf einen
Reiz ist von dessen jeweiliger emotionaler Bewertung abhängig.
Dementsprechend folgen spezifische biochemische Prozesse auf
verschiedene Affekte.}







\subsubsection{Diathese-Stress-Modell bzw. transaktionales Stressmodell}

Unter \EE{Diathese} versteht man die angeborene bzw.
erworbene Verletzlichkeit (Vulnerabilität) gegenüber
Krankheitsauslösern, oder im weiteren Sinne auch Stressoren, d.\,h. eine
interindividuelle Grundausstattung mit der Personen Stress
begegnen und somit anfälliger oder weniger anfällig für
Stressfolgen sind.\ZFootnote{(Davison \& Neale, 1988)}

Als \T{Coping} bezeichnet man ein Verhalten, um eine
stressauslösende Situation zu
bewältigen.\ZFootnote{\Lang{engl.} \Speech{to cope with}: 
fertig werden mit}.

\Commentary[Diathese-Stress-Modell bzw. transaktionales
Stressmodell]{Unter \EE{Diathese} wird die angeborene bzw. erworbene
Verletzlichkeit (Vulnerabilität) gegenüber
Krankheitsauslösern, oder im weiteren Sinne auch Stressoren, verstanden, d.\,h. eine
interindividuelle Grundausstattung mit der Personen Stress
begegnen und somit anfälliger oder weniger anfällig für
Stressfolgen sind. [Davison \& Neale, 1988]

Davon ausgehend erweitert das transaktionale Stressmodell
[Lazarus \& Launier, 1981] diese Grundausstattung um den
Begriff \I{Coping}, worunter Verhalten verstanden wird, um eine stressauslösende Situation zu bewältigen. D.\,h. nachdem
ein Individuum eine Bedrohung/Gefährdung antizipiert hat,
bewertet es diese und versucht gegebenenfalls mit ihr fertig
zu werden (englisch to cope with = fertig werden mit).
  }

\begin{itemize}
  \item \EE{Primäre Bewertung}: \Speech{\enquote{Ist der
  Reiz für mich überhaupt relevant?}} Wenn ja:
  \Speech{\enquote{Ist er eine Herausforderung oder eine
  Bedrohung?}} Die Beantwortung dieser Fragen hängt von der Einschätzung des eigenen Könnens  sowie vom
  Aktivationsgrad ab.
  
  \item \EE{Sekundäre Bewertung}: \Speech{\enquote{Habe ich
  Handlungsmöglichkeiten?}} Wird diese Frage mit
  \Speech{\enquote{Nein}} beantwortet, entsteht Stress; wird
  sie mit \Speech{\enquote{Ja}} beantwortet, kann das
  Individuum entsprechend reagieren und geht zum Coping
  über.
  
  \item \EE{Coping}: problemorientiert
  
  Ziel: Veränderung
  des Stressors durch entsprechendes Problembewältigungsverhalten
  oder
  gedankliche oder emotionale Neubewertung
  
  \item \EE{Neubewertung} der Situation.
\end{itemize}
\A{ZIMBARDO!!!}


\Fazit{Das Diathese-Stressmodell  postuliert eine
interindividuelle tw. angeborene, tw. sozialisierte
Anfälligkeit für Stressgeschehen. Im transaktionalen
Stressmodell (Coping-Modell) schätzt eine Person
entsprechend dieser persönlichen Grundausstattung in einer
primären Bewertung die Relevanz des Reizes ein, in einer
sekundären Bewertung seine persönlichen Möglichkeiten, den
Stressor zu bewältigen und geht ggfs. zum Coping
über.}

\Commentary[Theorien Stressmodelle]{Den vorgestellten Theorien entsprechend setzt sich das
Phänomen \enquote{Stress} also aus folgenden Komponenten zusammen:
dem ideellen \EE{Gleichgewicht (Homöostase)} einer Person,
das, abhängig von ihrer \EE{individuellen angeborenen bzw.
erworbenen Stressvulnerabilität, persönlichen
Bewertungsprozessen} zur Relevanz und der zur Verfügung
stehenden Handlungsmöglichkeiten, von \EE{einem oder
mehreren Reizen bedroht} wird \EE{(=Stressoren)}.
\EE{Allostatische Regelkreise} initiieren in mehrphasigem
Ablauf etwaige \EE{Copingversuche} auf
gedanklich-emotionaler oder Verhaltensebene, erneute
Bewertungsprozesse schätzen den Erfolg oder Misserfolg ein.
Intensive kritische oder chronisch belastende Situationen
können \EE{Fehlregulationen} in Gang setzten, die
pathologische Folgen haben können (siehe Kapitel
\enquote{Stressauswirkungen}).

Hervorgehoben werden soll, dass demzufolge die
Bewertungsprozesse und ein daran geknüpftes Gefühl der Handlungsmöglichkeit einen
maßgeblichen Einfluss auf das Stresserleben ausübt. Daraus
abgeleitete Implikationen zur positiven Stressbewältigung
sollen in den Kapiteln \enquote{Stressbewältigung} und \enquote{Grundsätze
der Stressvermeidung} behandelt werden.}



\subsection{Aktivierung, Aktivierungsniveau, Überforderung, Unterforderung}

\label{topic:ueberforderung-psi}
\label{topic:unterforderung-psi}



Unter \Z{Aktivierung} ist ein \EE{allgemeiner nervöser
Erregungszustand} des Organismus zu verstehen, der in
inneren Repräsentationen (Emotionen) bzw.
psychophsyiologischen Veränderungen seinen Ausdruck findet.
Er kann in Abhängigkeit von Bewusstseinszuständen (Wachheit,
Müdigkeit, Schlaf), Gefühlslagen (Entspannung, Anspannung,
Nervosität etc.) und Umweltbedingungen (stressarme vs.
stressreiche Umwelt) variieren. Indikatoren für Aktivierung
können sein:
\begin{itemize}
  \item Erhöhung der Herz-, Atem- bzw. Lidschlagfrequenz
  \item Anstieg des Fingerpulsvolumen, Hautwiderstand sinkt 
  wegen Anstiegs der Hautleitfähigkeit durch erhöhte
  Schweißproduktion
  \item Anstieg des Muskeltonus
  \item Gefühle der Anspannung/Aufregung
\end{itemize}


\EE{Aktivierung und Leistung stehen
miteinander in Beziehung}\ZFootnote{Yerkes-Dodson-Gesetz \cite{Yerkes1908}} im Sinne eines \E{umgekehrt
U-formiger Zusammenhangs}, d.\,h. bei sehr
geringerem und sehr hohem Aktivierungsniveau ist die resultierende Leistung niedrig, bei mittlerer
Aktivierung ist der Leistungsresultat am höchsten. 
Daraus
lässt sich ableiten, dass für leichte Aufgaben ein höheres
Aktivierungsniveau günstiger ist als für schwere Aufgaben --
deshalb lassen sich z.\,B. subjektiv leichte Aufgaben,
z.\,B. Routine-Aufgaben, vor Publikum (das in der Regel eine
höhere Aktivierung hervorruft) besser bewältigen als
schwierige Aufgaben.

\Commentary[Formen der Aktivierung]{
Es werden drei spezielle Formen der Aktivierung
unterschieden: Orientierungsreaktion, Defensivreaktion und
Schreckreflex: 

Die \EE{Orientierungsreaktion} ist eine
unspezifische Aktivierung bei Veränderungen in der
Reizumgebung im Sinne eines
\Speech{\enquote{Was-ist-das?}}-Reflexes. Sie ist besonders
ausgeprägt bei biologischer preparedness (eine angeborene Bereitschaft) oder gelernten bzw.
konditionierten Reizen. Eine Orientierungsreaktion tritt
auf, wenn ein neuer Stimulus im Wahrnehmungsbereich
erscheint. Ihre Intensität ist dabei proportional zur
Neuheit des Reizes, d.\,h. sie fällt umso stärker aus, je mehr
der Reiz von Bekanntem abweicht und wird auch begleitet von
charakteristischen Veränderungen in verschiedenen
Funktionssystemen des Körpers, z.\,B.
\\Sinnesorgane: die Reizschwellen sinken.
\\Haut: der Hautwiderstand sinkt.
\\Muskulatur: Anspannung, Aktivierung notwendiger Muskeln.
\\Blutgefäße: periphere Vasokonstriktion, Vasodilatation der
ZNS-, Lungen-, Herz-Blutgefäße.
\\Herz: zunächst sinkt die Herzfrequenz aufgrund peripherer
Vasokonstriktion, anschließend steigt sie.

Bei mehrmaliger Darbietung desselben (neuen) Reizes erfolgt
eine Gewöhnung (Habituation), d.\,h. die Intensität der
Orientierungsreaktion nimmt ab. Änderungen des Reizes führen
zu Dishabituation, es tritt wieder die ursprüngliche
Orientierungsreaktion auf.

Die \EE{Defensivreaktion} ist eine Reaktion zur Verteidigung
gegenüber schmerzhaften bzw. aversiven Reizen: das
Individuum versucht den Reizkontakt abzubrechen und leitet
Abwehrmaßnahmen ein. Diese Reaktion unterliegt keiner
Habituation.

Unter dem \EE{Schreckreflex} wird eine spontane, willentlich
nicht kontrollierbare Reaktion verstanden. Schreckreflex
unterbricht die Informationsaufnahme und motorische
Aktivität bei neuen und aggressiven Reizen, um Verhalten und
Aufmerksamkeit rasch auf potentiell gefährliche Reize
umzuorientieren, z.\,B. Blinzeln bei plötzlich auftretendem
Lärm.
}



Ein gewisses Aktivierungsniveau ist
erforderlich, um Leistungen überhaupt erbringen zu können --
zu hohe und zu niedrige Aktivierung schwächt die
Leistungsfähigkeit jedoch ab (Yerkes-Dodson-Gesetz).
Aktivierung wird begleitet von typischen
psychophysiologischen Veränderungen. Unter der
Orientierungsreaktion ist eine aktivierende Reaktion bei
Auftreten neuer Reize, weitere aktivierende Reaktionen sind
die Defensivreaktion (Verteidigung gegen unerwünschte Reize)
und der Schreckreflex.


Je nach persönlicher Veranlagung, Sozialisation und
Habituation (s.\,o.) verfügt jeder Mensch über sein
\EE{individuelles Erregungsniveau bzw. Aktivierungsprofil},
welches sich entsprechend jeweiliger Lebensumstände und
Stimmungslage auch
verändern bzw. verschieben kann. Dieser Umstand kann
Konsequenzen in Bezug auf pathologische Entwicklungen sowie
das subjektive Gefühl von Über- bzw. Unterforderung haben.
% (siehe Abbildungen).













% 
% Die Abbildung verdeutlicht, dass Personen mit höherem
% persönlichen Aktivierungsniveau gegenüber pathologischen
% Veränderungen vulnerabler sein werden.



Ausgehend davon, dass mit einem erhöhtem Erregungsniveau
auch die Anforderungen einer Aufgabe subjektiv als
herausfordernder wahrgenommen werden, fühlen sich --
entsprechend der Abbildung -- \enquote{übererregte} Menschen
schneller überfordert, als solche mit durchschnittlichem
Aktivierungslevel.

\A{Keywords: Physiologie von Stress, Kritik an \enquote{dem Stress,
Eu- Distress}, evolutionärer Nutzen, Stressmodelle/
-theorien, kognitive Bewertung der Gegebenheiten einer
Situation, persönliche Ressourcen, (Aktivierung,
individuelles Aktivierungsprofil).}

Mit dem Begriff Aktivierung ist ein allgemeiner Erregungszustand gemeint, der seinen Ausdruck in unterschiedlichen Gefühlszuständen (z.B. Gefühl der Anspannung/Aufregung) und physiologischen Veränderungen (z.B. erhöhte(r) Herzschlag, Muskeltonus und Atemfrequenz) findet und von unterschiedlichen Bewusstseinszuständen und Umweltbedingungen beeinflusst wird. Zahlreiche psychologische Studien verweisen auf einen umgekehrt U-förmigen Zusammenhang zwischen Aktivierung und Leistungsfähigkeit hin: d.h. die optimale Leistung lässt sich in mittleren Erregungszuständen erbringen, zu niedrige (Unterforderung) oder zu hohe Aktivierung (Überforderung) schwächt die Leistungsfähigkeit und verursacht gleichzeitig auch Unbehagen (z.B. Anspannungsgefühl wegen Langeweile bzw. Überforderung, Ängstlichkeit, Verunsicherung). In der Regel suchen wir -sofern möglich, entsprechend des Gleichgewichtmodells - Situationen auf, die uns in ein angenehmes Aktivierungsniveau versetzen, d.h. bei zu niedriger Aktivierung gehen wir Aktivitäten/Beschäftigungen nach (Gespräch, Hobby, Sport etc.), bei hoher Aktivierung ziehen wir uns zurück, versuchen zu entspannen.

Je nach persönlicher Veranlagung und Lebenserfahrung verfügt jeder Mensch über sein individuelles Erregungsniveau (z.B. Menschen die generell eine höhere oder niedrigere Grunderregung aufweisen)und Aktivierungsprofil (z.B. Menschen die sensibler bzw. weniger sensibel auf Umwelteinflüsse reagieren). Beides kann sich entsprechend jeweiliger Situationen/Lebensumstände (z.B. Versorgung einer Schürfwunde vs. Reanimation)   und
Stimmungslage (z.B. unangenehme negative Stimmungslage erhöht Anspannungslevel und Stresssensibilität)auch verändern bzw. verschieben kann. Dies kann das subjektive Gefühl von Über- bzw. Unterforderung beeinflussen und pathologische Entwicklungen begünstigen oder vor solchen schützen.



Ableitungen für den Sanitätsdienst: 

- Je nach persönlichem Aktivierungsprofil, Tagesverfassung und Situation liegen bei den Teammitgliedern eines Sanitätsteams unterschiedliche Aktivierungsniveaus vor. Diesem Umstand sollte mit Verständnis und Toleranz Rechnung getragen werden. Ähnliche Situationen können an einem Tag als gut bewältigbar, am nächsten als überfordernd erlebt werden – ebenso verhält es sich von Person zu Person unterschiedlich.
- Im Sinne der Stressvermeidung wäre es wichtig, Sensibilität für sein persönliches Erregungsprofil zu entwickeln, z.B. „Fühle ich mich generell eher angespannt oder aufgeregt? Wie fühle ich mich jetzt gerade? Auf welche Auslöser/Situationen reagiere ich besonders stark?“
- Es bestehen Möglichkeiten der Einflussnahme, Techniken sein persönliches Aktivierungsniveau zu beeinflussen können erlernt/gefördert werden: Aktivierung durch persönliche Zielsetzung, betrachten einer Aufgabe als Herausforderung etc. ebenso wie Entspannung durch pos. Selbstinstruktion („ich schaffe das“), Atemtechniken (bewusstes ruhiges Atmen), bewusste Pausen etc.



\section{Stressauswirkungen, Früherkennung, Beanspruchungsfolgen}



\subsection{Stressauswirkungen}
\label{topic:stressauswirkungen-psi}

Prinzipiell können sich die Auswirkungen von Stress in allen
Facetten des Erlebens und Verhaltens von Betroffenen
wiederspiegeln, d.\,h. global betrachtet im \EE{persönlichen
Wohlbefinden} (z.\,B. unspezifische Gefühle des
Unwohlseins), \EE{auf emotionaler Ebene} (z.\,B. gedrückte
Stimmung, Frustrationsgefühle, Gefühle der
Unsicherheit/Ängstlichkeit), \EE{auf gedanklicher Ebene}
(z.\,B. neg. Gedanken, Grübeln, Selbst- oder
Fremdentwertung), \EE{im Antrieb} (z.\,B. Energielosigkeit,
Antriebsschwächen), \EE{in der Motivation} (z.\,B. Unlust,
Interessensverlust, Perspektivlosigkeit), in der
\EE{persönlichen Einstellung bzw. Wertesystemen} (z.\,B.
Zynismus, negativistische Weltsicht, Pessimismus), in der
\EE{zirkadianen Rhythmik} (Ein- bzw. Durchschlafstörungen,
Müdigkeit, Überdrehtsein) sowie \EE{im Umgang mit sich
selbst} (z.\,B. selbstschädigendes Verhalten durch
Alkohol/Drogen, über seine Grenzen gehen, Gefühle
unterdrücken/ignorieren) bzw. \EE{mit seiner sozialen
Umwelt} (z.\,B. soziale Normen brechen,
gereiztes/feindseliges Reagieren, sozialer Rückzug).
Desweiteren sei auf mögliche \EE{körperliche Warnsignale}
hingewiesen, wobei häufig in diesem Zusammenhang von
folgenden Beschwerden berichtet wird: erhöhte
Schmerzempfindlichkeit, Rücken-, Kopf- oder Bauchschmerzen,
erhöhter Blutdruck, Rastlosigkeit oder Sodbrennen.

\subsection{Früherkennung}
\label{topic:stressfrueherkennung-psi}

Die oben dargestellte, weit gefächerte Bandbreite an
möglichen Ebenen, auf denen sich die Folgen von Stress
manifestieren können, macht eine Früherkennung nicht
unbedingt leicht. Hervorzuheben ist, dass die \EE{eigene
Selbstwahrnehmung} (v.\,a. im Bezug auf Gefühle und
Stimmungen) und \EE{Introspektion} den erstmöglichen Zugang
bietet, da der Betroffene selbst auf gedanklich-emotionaler
Ebene wohl die erste Chance bekommt, die Vorboten von
Distress zu erkennen bzw. am eigenen Leib zu \enquote{spüren}. Da
viele der o.\,g. Auswirkungen aber sehr \EE{subtil bzw.
schleichend sukzessive} erfolgen bzw. auch durch unbewusste
Abwehrmechanismen (z.\,B. Verdrängung) sich der bewussten
Wahrnehmung häufig entziehen, kann eine \EE{Sensibilisierung
der Eigenwahrnehmung von Gefühlen} bzw. \EE{Schärfung der
Selbstreflexion durch Inter- und/oder Supervision} einen den
Selbstschutz fördern. Nehmen Betroffene die Warnsignale
nicht wahr, könnte eine ausgeprägte \EE{Offenheit für
Feedback} über das eigene Verhalten durch andere Personen,
z.\,B. Angehörige, Freunde, Kollegen, Supervisor, Peergroup
etc. (Fremdwahrnehmung) einen möglichen Zugang darstellen.
Diesbzgl. wäre eine \EE{entsprechende Teamkultur} eine
unbedingte Voraussetzung.


\subsection{Beanspruchungsfolgen}
\label{topic:beanspruchungsfolgen-psi}

\begin{itemize}
  \item \EE{Erhöhte Vulnerabilität}:
Generell kann Stress eine erhöhte Vulnerabilität
(=Verletzlichkeit) zur Folge haben: im sozialpsychologischen
Kontext mag damit eine erhöhte \EE{emotionale
Verletzlichkeit} gegenüber Belastungen, Enttäuschungen und
Frustrationen im zwischenmenschlichen Bereich sowie
eingeschränkte Handlungsmöglichkeiten, adäquat darauf
reagieren zu können, gemeint sein. Ebenso ist unter dem
Begriff die persönliche Anfälligkeit für \EE{gesundheitliche
Erkrankungen} gemeint sein, z.\,B. durch die oben
beschrieben Schwächung der Immunkompetenz.
\item 
\EE{Frustration}:
Ausgehend von der Prämisse, dass unbewältigter Stress
Frustration erzeugt, lassen sich zwei potentielle Richtungen
unterscheiden, in die sich diese kanalisieren kann:
Frustrationsgefühle können sich nach außen richten, durch
das gereizte Reagieren auf Mitmenschen (\EE{Aggression}: im
Sinne von erhöhter Reizbarkeit und geringerer
Frustrationstoleranz), oder nach innen, durch persönliche
Selbstentwertung (\EE{Depression}: im Sinne von neg.
Gedanken, Gefühlen der Unzulänglichkeit/Überforderung und
gedrückter Stimmung).

\item 
\EE{Zynismus}\ZFootnote{\Lang{griech.}
kynismós, wörtlich \enquote{Hündigkeit} im Sinne von
\enquote{Bissigkeit}, von \Lang{griech.} kyon,
\enquote{Hund}) bezeichnete ursprünglich die Lebensanschauung und -weise der antiken Kyniker
(philosophische Richtung der griechischen Antike, deren
Zentrale Grundhaltung der ethischen Skeptizismus darstellt.} bezeichnet zum einen eine
\EE{Haltung, die in (oft absichtlich) verletzender Weise
die Wertvorstellungen anderer herabsetzt oder missachtet} und
zum anderen auch eine \EE{Haltung, die moralische Werte
grundsätzlich in Frage stellt und sich darüber hinaus
manchmal auch über sie lustig macht}. Auch eine hieraus
folgende berechnend-amoralische Einstellung und
Verhaltensweise kann Ausdruck dieser Haltung sein. Zynismus
kann Folge und Anzeichen von Resignation sein. 

Dies legt ein erhöhtes Risiko für Angehörige der Gesundheitsberufe nahe,
da die Konfrontation mit hoffnungslos erscheinenden
Situationen ein Gefühl der Hilflosigkeit erzeugen kann.\ZFootnote{Dass
eine zynische Grundhaltung nicht nur im zwischenmenschlichen
Bereich Probleme verursachen (z.\,B. durch neg. Bewertung
durch Menschen aufgrund  wahrgenommener Feindseligkeit im
Umgang miteinander), sondern auch mit ungünstigen
gesundheitlichen Konsequenzen einhergehen kann, zeigt die
folgende Untersuchung. Die Studie \enquote{Hostility, the Metabolic
Syndrome, and Incident Coronary Heart Disease} \cite{Niaura2002} ist
dem Zusammenhang zwischen Feindseligkeit und der Schwere
koronarer Herzerkrankungen nachgegangen. Zwei potentielle
Wirkmechanismen wurden für den hohen positiven Zusammenhang
angenommen, nämlich, dass Personen mit hohen Werten in dem
Merkmal Feindseligkeit möglicherweise aufgrund ihrer
ablehnend-zynischen Einstellung von anderen Menschen weniger
soziale Unterstützung erfahren bzw. mit einem generell
ungünstigen psychosozialen Risikoprofil einhergehen, das
durch mangelnde soziale Unterstützung, häufige
Ärgersituationen sowie fehlende bewältigungsregulative
Ressourcen gekennzeichnet ist.}
\end{itemize}

\subsection{Pathologische Phänomene}

Im Extremfall können andauernde oder sehr intensive
Belastungssituationen klinische Zustandsbilder im Sinne
psychischer Erkrankungen auslösen, wobei betont sei, dass er
Übergang von \enquote{psychisch gesund} zu \enquote{psychisch krank} ein
fließender ist \A{(siehe Exkurs \enquote{Diagnosestellung})}.
Exemplarisch sei hier auf die Phänomene \enquote{Burnout} und
\enquote{Traumafolgestörungen} eingegangen, da diese Folgen
chronischen bzw. besonders intensiven emotionalen
Belastungsgeschehens darstellen können.

\A{Exkurs: Diagnosestellung}

\subsubsection{Burnout}

\label{topic:burn-out}

Eine genaue Definitionen des
Burnout-Syndroms im Sinne einer Diagnose nach
Klassifikationssystem ist schwierig.\ZFootnote{Burnout:
Geprägt wurde der Begriff 1974 durch Psychoanalytiker Freudenberger. Die Abgrenzung ist deshalb besonders
schwierig, da es sich dzt. noch um ein beforschtes
Konstrukt handelt und die differentialdiagnostische Abgrenzung zu ähnlichen,
psychischen Erkrankungen (v.\,a. depressiven Erkrankungen)
sich als entsprechend schwierig erweist, weswegen bisher
nach ICD-10 \cite{ICD10-De-2011}
ausschließlich die \EE{Hilfsdiagnose Z73.0
\enquote{Ausgebranntsein}}, \EE{\enquote{Zustand der totalen
Erschöpfung}} existiert.} Im Folgenden werden zwei
unterschiedliche Abgrenzungsversuche dargestellt:

\begin{itemize}
  \item \EE{Persönlichkeitsorientierte Ansätze}: Sie stellen
  Persönlichkeit des Helfenden in den Vordergrund und sehen
  die Ursache für Burnout in einer \EE{Diskrepanz zwischen dem
  Helferideal} -- also den eigenen Ansprüchen als Helfer --
  \EE{und der Wirklichkeit des Helfens}.\ZFootnote{\Speech{\enquote{Burnout wird in Gang gesetzt durch Autonomieeinbußen in gestörten Auseinandersetzungen
  des Individuums mit seiner Umwelt, genauer: durch die innere
  Repräsentation solcher Interaktionen als gestörter und das
  Scheitern bei ihrer Bewältigung}} [Burisch, 1994]}
  
  \Speech{\enquote{Before you get burned-out you have to
  get on fire}}
  \begin{itemize}
    \item Perfektionismus, Zwanghaftigkeit
    \item Hohes Kontrollbedürfnis
    \item Übermäßiges Verantwortungsbewusstsein
    \item Hang zu übermäßigen und unrealistischen Schuldgefühlen
    \item Idealismus, Übereifer
    \item Verkopfung bzw. Unterdrückung von Gefühlen
    \item Wunsch alles selber zu machen bzw. kein Erbitten von Hilfe
    \item Schwierigkeit freie Zeit zu nehmen und zu genießen
  \end{itemize}
  \item \EE{Organistationsorientierte Ansätze}: Diese betonen die
  Relevanz von sozialen, Arbeits- und
  Organisations-Bedingungen und sehen die Entstehungsgründe
  für Burnout im emotional beanspruchenden und
  erschöpfenden Umgang mit Menschen.\ZFootnote{\Speech{\enquote{ Burnout ist
  ein Syndrom emotionaler Erschöpfung, Depersonalisierung und
  reduzierter persönlicher Leistungsfähigkeit, das bei
  Individuen, die in irgendeiner Weise mit Menschen arbeiten,
  auftreten kann}} [Maslach \& Jackson, 1984]}\ZFootnote{Maslach, 1982;
  Frieling \& Sonntag, 1999; Bergner, 2003}
  \begin{itemize}
    \item Sozialer Stress: Interaktion mit Patienten,
    Überforderung im Umgang mit
    Krankheiten, hohe Verantwortung, Mangel an positivem
    Feedback, Hierarchieprobleme, schlechte Teamarbeit, Druck
    von Vorgesetzten
    \item Unrealistische Vorstellung der
    Patienten/Gesellschaft: Arzt/Helfer ist omnipotent, darf
    keine eigenen Bedürfnisse haben, Normen eines guten
    Helferverhaltens \item Arbeitsbedingungen: ungünstige
    Arbeitsorganisation, personelle Unterbesetzung, hoher
    Zeitdruck, kaum Möglichkeit zur Arbeitsplatzgestaltung
    \item Mehr verwaltungstechnische und
    betriebswirtschaftliche Aufgaben
    \item Abnahme der Autonomie der Ärzte/Helfer durch
    stärkere Gesetzgebung im Gesundheitswesen
    \item Verminderung der immateriellen Entschädigung
    \item Verminderung der materiellen Entschädigung
  \end{itemize}
\end{itemize}



\TextSlide{\TxSymptome}
{Burnout kann sich durch eine Reihe von höchst unterschiedlichen Symptomen bemerkbar machen. 
Man kann unterscheiden\ZFootnote{nach Schaufeli (1992), zitiert
nach \cite{BurischBurnOut3}}:

\begin{itemize}
  \item \EE{Psychische Symptome}
  emotional: großer Widerstand täglich zur Arbeit zu gehen,
  Gefühle des Versagens, Ärgers und Widerwillens,
  Schuldgefühle, Entmutigung und Gleichgültigkeit, Misstrauen
  und paranoide Vorstellungen, Frustration. kognitiv:
  Rigidität im Denken und Widerstand gegen Veränderungen,
  Projektionen, Konzentrationsstörungen. motorisch: nervöse
  Ticks, Verspannungen.
  \item \EE{Physische Symptome} und
  psychosomatische Beschwerden: tägliche Gefühle von Müdigkeit
  und Erschöpfung, große Müdigkeit nach dem Arbeiten,
  Schlafstörungen, sexuelle Probleme. Erkrankungen: häufige
  Erkältungen und Grippe, häufige Kopfschmerzen,
  Magen-Darm-Beschwerden. Physiologische Reaktionen: erhöhter
  Herzschlag, erhöhter Pulsfrequenz, erhöhter Cholesterinspiegel.
  \item \EE{Symptome auf der Verhaltensebene}:
  Individuelle Verhaltensweisen: exzessiver Drogengebrauch,
  Tabakgenuss, Alkoholkonsum und/oder Kaffeekonsum, erhöhte
  Aggressivität. Verhalten in der Arbeit: häufiges Fehlen am
  Arbeitsplatz, längere Pausen, verminderte Effizienz.
  \item \EE{Soziale Symptome}:
  Im Umgang mit Klienten: Verlust von positiven Gefühlen den
  Klienten gegenüber, Verschieben von Klientenkontakten,
  Widerstand gegen Anrufe und Besuche von Klienten,
  Unfähigkeit sich auf Klienten zu konzentrieren oder ihnen
  zuzuhören. Im Umgang mit Kollegen: Isolierung und Rückzug,
  Vermeidung von Arbeitsdiskussionen mit Kollegen. außerhalb
  der Arbeit: Ehe- und Familienprobleme, Einsamkeit.
  
  \item \EE{Problematische Einstellung}:
  Im Umgang mit Klienten: Stereotypisierung von Klienten,
  Zynismus, schwarzer Humor, verminderte Empathie,
  Demonstration von Machtlosigkeit. In der Arbeit/
  Einrichtung: negative Arbeitseinstellung, Desillusionierung,
  Verlust von Idealismus.
  \item \EE{Warnsymptome}: In der Entstehungsphase eines Burnout-Syndroms gibt es
  charakteristische Merkmale beispielsweise:
  \begin{itemize}
    \item extremes Engagement für ein bestimmte Ziel/bestimmte Ziele.
    \item Hyperaktivität
    \item chronische Müdigkeit und körperliche sowie geistige Erschöpfung
    \item praktisch pausenloses arbeiten und das Arbeiten zum
    wichtigsten Lebensinhalt machen
    \item Verzicht auf Urlaub, Erholungsphasen und Entspannungsphasen
    \item sich unentbehrlich und vollkommen fühlen
    \item Abwertung anderer, um sich selbst (künstlich) zu erhöhen
    \item sozial unangemessenes Verhalten
    \item Zunehmendes Ignorieren eigener Bedürfnisse
    \item möglicherweise Gewichtsabnahme und
    Mangelerscheinungen durch unzureichende Ernährung bei
    gleichzeitig erhöhtem Bedarf
    \item Verdrängen von Misserfolgen, bei gleichzeitigem
    verantwortlich machen anderer
    \item Beschränkung von sozialen Kontakten auf den
    Arbeitsbereich – private Kontakte werden vermieden, auch
    zum Lebenspartner und gegebenenfalls Kindern und anderen
    Verwandten
    \item selbstschädigendes Verhalten/Ablenkung durch
    Alkohol, Suchtmittel, Glückspiel, Internet-
    /Computeraktivitäten und Sexualität.
    \item Konzentrationsschwäche
    \item Schlafstörungen (Einschlafstörungen und Durchschlafstörungen)
    \item Drehschwindel und Neigung zu Tinnitus
    \item Angstzustände
  \end{itemize}
\end{itemize}
}
{
\begin{itemize}
  \item {Psychische Symptome}
  \item {Physische Symptome und psychosomatische Beschwerden}
  \item {Symptome auf der Verhaltensebene}
  \item {Soziale Symptome}
  \item {Problematische Einstellung}
  \item {Warnsymptome}
\end{itemize}
}
{}
{}
{}

\begin{itemize}
  \item \EE{Psychische Symptome}
  emotional: großer Widerstand täglich zur Arbeit zu gehen,
  Gefühle des Versagens, Angstzustände, Ärgers und Widerwillens,
  Schuldgefühle, Entmutigung und Gleichgültigkeit, Misstrauen
  und paranoide Vorstellungen, Frustration. kognitiv:
  Eingeschränkte Eigenwahrnehmung, Zunehmendes Ignorieren
  eigener Bedürfnisse
  Rigidität im Denken und Widerstand gegen Veränderungen,
  Projektionen, Konzentrationsstörungen. motorisch: nervöse
  Ticks, Verspannungen, Hyperaktivität.
  \item \EE{Physische Symptome} und
  psychosomatische Beschwerden: tägliche Gefühle von Müdigkeit
  und Erschöpfung, große Müdigkeit nach dem Arbeiten, körperliche sowie geistige Erschöpfung,
  Schlafstörungen, sexuelle Probleme. Erkrankungen: häufige
  Erkältungen und Grippe, häufige Kopfschmerzen,
  Magen-Darm-Beschwerden, Drehschwindel und Tinnitus. Physiologische Reaktionen: erhöhter
  Herzschlag, erhöhter Pulsfrequenz, erhöhter Cholesterinspiegel, Gewichtsabnahme und
    Mangelerscheinungen durch unzureichende Ernährung bei
    gleichzeitig erhöhtem Bedarf
.
  \item \EE{Symptome auf der Verhaltensebene}:
  Individuelle Verhaltensweisen: 
  selbstschädigendes Verhalten/Ablenkung durch Konsum von Tabak, Kaffee 
    Alkohol, Suchtmittel, Glückspiel, Internet-
    /Computeraktivitäten und Sexualität. erhöhte
  Aggressivität. Verhalten in der Arbeit: häufiges Fehlen am
  Arbeitsplatz, längere Pausen, verminderte Effizienz.
  \item \EE{Soziale Symptome}:
  Im Umgang mit Klienten: Verlust von positiven Gefühlen den
  Klienten gegenüber, Verschieben von Klientenkontakten,
  Widerstand gegen Anrufe und Besuche von Klienten,
  Unfähigkeit sich auf Klienten zu konzentrieren oder ihnen
  zuzuhören. Im Umgang mit Kollegen: Isolierung und Rückzug,
  Vermeidung von Arbeitsdiskussionen mit Kollegen. Außerhalb
  der Arbeit: Beschränkung von sozialen Kontakten auf den
  Arbeitsbereich – private Kontakte werden vermieden Ehe-
  und Familienprobleme, Einsamkeit.
  
  \item \EE{Problematische Einstellung}:
  Im Umgang mit Klienten: Stereotypisierung von Klienten,
  Zynismus, schwarzer Humor, verminderte Empathie,
  Demonstration von Machtlosigkeit. In der Arbeit/
  Einrichtung: sozial unangemessenes Verhalten, Verdrängen
  von Misserfolgen, bei gleichzeitigem verantwortlich machen
  anderer, negative Arbeitseinstellung, Desillusionierung,
  Verlust von Idealismus, Abwertung anderer, um sich selbst
  (künstlich) zu erhöhen, sich unentbehrlich und vollkommen
  fühlen, Verzicht auf Urlaub, Erholungsphasen und
  Entspannungsphasen, praktisch pausenloses arbeiten und das
  Arbeiten zum wichtigsten Lebensinhalt machen, extremes
  Engagement für ein bestimmte Ziel/bestimmte Ziele.
\end{itemize}

     


\subsubsection{Traumafolgestörungen}

% \EE{Ätiologie}

\Text{\TxBeschreibung}
{Eine wesentliche Grundvoraussetzung für die Entstehung von
Traumafolgestörungen ist eine (oder mehrere) vorangegangene
punktuelle oder andauernde Konfrontation mit einem
außergewöhnlich intensiven emotionalen
\EE{Belastungserlebnis} (\E{=traumatisches Erlebnis,
psychisches Trauma = griech. für Wunde}), auf das die Person \EE{mit intensiver Furcht,
Hilflosigkeit oder Entsetzen} reagiert. 
Für eine
exemplarische Auflistung möglicher Traumaereignisse siehe
Fehlregulationsmodell. 
Hervorzuheben ist jedoch deutlich,
dass  unabhängig von der Art des Ereignisses das
\E{subjektive Erleben des/r Betroffenen entscheidend} ist,
also \Speech{\enquote{nicht die objektive, sondern die subjektive
Bedrohlichkeit eines Traumas die posttraumatische
Symptombelastung vorhersagt}}\footnote{Kilpatrick, Best,
Veronen, Villeponteaux \& Amick-Mc-Mullen, 1986}. Um ein
ungefähres Bild darüber zu bekommen, wie etwaige Folgen einer
Traumatisierung sich gestalten können, sind im Folgenden
Abschnitt die jeweiligen Diagnosebeschreibungen laut ICD-10
und phänomenolgischen Leitsymptome dargestellt
\cite{ICD10-De-2011}.}
{}
{}
{}
{}




% \EE{Phänomenologie}

\Text{Akute Belastungsreaktion (F43.0)}
{Eine vorübergehende Störung, die sich bei einem psychisch
nicht manifest gestörten Menschen als Reaktion auf eine
außergewöhnliche physische oder psychische Belastung
entwickelt, und die im allgemeinen innerhalb von Stunden
oder Tagen abklingt. Die individuelle Vulnerabilität und die
zur Verfügung stehenden Bewältigungsmechanismen
(Coping-Strategien) spielen bei Auftreten und Schweregrad
der akuten Belastungsreaktionen eine Rolle. Die Symptomatik
zeigt typischerweise ein gemischtes und wechselndes Bild,
beginnend mit einer Art von \enquote*{Betäubung}, mit einer
gewissen Bewusstseinseinengung und eingeschränkten Aufmerksamkeit,
einer Unfähigkeit, Reize zu verarbeiten und
Desorientiertheit. Diesem Zustand kann ein weiteres
Sichzurückziehen aus der Umweltsituation folgen (bis hin zu
dissoziativem Stupor, F44.2) oder aber ein
Unruhezustand und Überaktivität (wie Fluchtreaktion oder
Fugue). 

Vegetative Zeichen panischer Angst wie Tachykardie,
Schwitzen und Erröten treten zumeist auf. Die Symptome
erscheinen im allgemeinen innerhalb von Minuten nach dem
belastenden Ereignis und gehen innerhalb von zwei oder drei
Tagen, oft innerhalb von Stunden zurück. Teilweise oder
vollständige Amnesie (F44.0) bezüglich dieser Episode
kann vorkommen. Wenn die Symptome andauern, sollte eine
Änderung der Diagnose in Erwägung gezogen werden.

Synonyme: Belastungsreaktion, Krisenreaktion, Psychischer Schock, Krisenzustand.

\subparagraph{Leitsymptome}
\begin{itemize}
  \item (emotionale) Taubheit
  \item Bewusstseinseinengung: eingeschränkte Aufmerksamkeit, \item beeinträchtigte Verarbeitung von Reizen
  \item Desorientierung
  \item Hyperaktivität, Unruhe, Angstzustände
  \item Rückzugsverhalten
  \item Rascher Eintritt der Symptome, Abklingen nach zwei bis drei Tagen
\end{itemize}
\subparagraph{Folgerungen für den Sanitätsdienst} Die
Wahrscheinlichkeit mit einem solchen oder ähnlichen Zustand
bei zu Versorgenden und/oder deren Angehörigen zu stoßen,
ist auf Patientenseite v.\,a. bei Unfallopfern,
Großschadensereignissen oder plötzlichen akuten
schweren oder lebensbedrohlichen Erkrankungen hoch einzuschätzen.
Entsprechend der  Symptomatik könnten folgende
Verhaltensregeln von Vorteil sein: 
\begin{itemize}
  \item einfache, kurze
  Anweisungen/Erklärungen,
  \item strukturierende Gesprächsführung,
  \item ruhiger Tonfall und beruhigender Umgang,
  \item Vermitteln von
  Sicherheit und Orientierungshilfen anbieten,
  \item Aufmerksamkeit
  aufs Hier und Jetzt lenken
\end{itemize}
-- alles nach dem Motto: \Speech{\enquote{Begleiten,
aber nicht aufdrängen}}, denn es handelt sich um eine \EE{normale
Reaktion} auf eine \EE{unnormale Situation} und in der Regel
verfügen die Betroffenen über \E{Selbstregulationsmechanismen}
zur Bewältigung. Ggfs. kann auch auf weitere
professionelle Hilfestellung (Krisenintervention,
Psychosozialer Dienst, Klinischer Psychologie,
Psychotherapie, telefonische Seelsorge etc.) verwiesen werden.

Analog verhält es sich mit dem Umgang, sollten Sanitäter
selbst davon betroffen sein, wobei hier etwaige Anzeichen
wahrscheinlich eher nach Ablauf des Einsatzes verstärkt
wahrgenommen werden. Nachbesprechungen (Debriefings),
Intervision oder \I{Supervision} können diesbzgl.
unterstützenden/präventiven Charakter haben und die
Selbstregulationsprozesse stützen. Ebenso stehen o.g.
Einrichtungen der professionellen Hilfe zur Verfügung,
sollte sich bei der/dem Betroffenen keine Erleichterung
einstellen. Eine entsprechende Teamkultur, diese
Hilfestellungen auch selbstverständlich zu nutzen, wäre
günstig.
}
{}
{}
{}
{}


\Text{Posttraumatische Belastungsstörung (F43.1)}
{\Z{\Lang{Syn.} \I[Stresssyndrom!postraumatisches]{Postraumatisches
Stresssyndrom} (\I{PTSS}), \Lang{engl.} \I{Post Traumatic
Stress Disorder} (\I{PTSD})}

\index{Belastungsstörung!posttraumatische}
\label{topic:ptss}
Diese entsteht als eine protrahierte bzw. \EE{verzögerte 
Reaktion auf ein belastendes Ereignis} oder eine Situation
kürzerer oder längerer Dauer, mit außergewöhnlicher
Bedrohung oder katastrophenartigem Ausmaß, die bei fast
jedem eine tiefe Verzweiflung hervorrufen würde.
Prädisponierende Faktoren wie bestimmte, z.\,B. zwanghafte
oder asthenische Persönlichkeitszüge oder neurotische
Krankheiten in der Vorgeschichte können die Schwelle für die
Entwicklung dieses Syndroms senken und seinen Verlauf
erschweren, aber die letztgenannten Faktoren sind weder
notwendig noch ausreichend, um das Auftreten der Störung zu
erklären. 

Typische Merkmale sind das \EE{wiederholte Erleben des
Traumas} in sich aufdrängenden Erinnerungen
(Nachhallerinnerungen, Flashbacks), Träumen oder \I[Albtraum]{Albträumen},
die vor dem Hintergrund eines andauernden Gefühls von
\EE{Betäubtsein} und \E{emotionaler Stumpfheit} auftreten. Ferner
finden sich Gleichgültigkeit gegenüber anderen Menschen,
Teilnahmslosigkeit der Umgebung gegenüber, Freudlosigkeit
sowie Vermeidung von Aktivitäten und Situationen, die
Erinnerungen an das Trauma wachrufen könnten. Meist tritt
ein Zustand von vegetativer Übererregtheit mit
Vigilanzsteigerung, einer übermäßigen Schreckhaftigkeit und
Schlafstörung auf. Angst und Depression sind häufig mit den
genannten Symptomen und Merkmalen assoziiert und
Suizidgedanken sind nicht selten. 

Der Beginn folgt dem
Trauma mit einer \EE{Latenz}, die wenige Wochen bis Monate dauern
kann. Der Verlauf ist wechselhaft, in der Mehrzahl der Fälle
kann jedoch eine Heilung erwartet werden. In wenigen Fällen
nimmt die Störung über viele Jahre einen chronischen Verlauf
und geht dann in eine andauernde Persönlichkeitsänderung
(F62.0) über.

\subparagraph{Leitsymptome}
\begin{itemize}
  \item realitätsnahes, intrusives \EE{Wiedererleben} (im
  Schlaf als Alpträume oder wach als Flashbacks/Erinnerungsbilder)
  -- evtl. ausgelöst durch traumaassoziierte Reize.
  \item \EE{Vermeidungsverhalten} gegenüber traumarelevanten
  Reizen, z.\,B. Gefühle, Gedanken, Gespräche, Aktivitäten,
  Personen oder Situationen
  \item \EE{emotionale Taubheit}, Entfremdung von anderen
  Menschen
  \item Hypervigilanz: erhöhte Reizbarkeit sowie
  übertriebene Schreckhaftigkeit
  \item Veränderte kognitiver
  Schemata/Kontrollüberzeugungen, z.\,B. verstärktes
  Misstrauen gegenüber Mitmenschen, Gefühl einer
  bedrohlichen Welt/Zukunft.
  \item Diagnosestellung erst nach vier Wochen ab dem
  Zeitpunkt der Traumatisierung bei anhaltender Symptomatik
  möglich.
\end{itemize}

\subparagraph{Folgerungen für den Sanitätsdienst} Der Umgang mit
PTSD-Betroffenen sollte stets ein umsichtiger und
reflektierter sein, da die Gefahr einer Retraumatisierung
(siehe unten) besteht. Vermeidungsverhalten sollte, sofern
es die Versorgungssituation erlaubt, möglichst respektiert
werden (z.\,B. im Anamnesegespräch nicht übermäßig
nachbohren, wenn der/die Betroffene(r) zu bestimmten Themen
schweigt oder lieber mit einer weiblichen einem männlichen
Sanitäter(in) sprechen möchte etc.). 

Sozial auffälliges
Verhalten kann auf mögliche Entfremdung bzw. erhöhtes
Misstrauen zurückzuführen sein, dennoch sind auch in der
empathischen Helferrolle die persönliche Grenzen klar
aufzuzeigen (siehe Sekundärtraumatisierung). Das Schaffen
einer ruhigen Atmosphäre kann Schreckreaktionen vermeiden,
jedenfalls könnten solche, sollten sie ausgelöst werden, als
\enquote{normale} Folge des Krankheitsbildes im Sinne der
Hypervigilanz bewertet werden. Hervorzuheben ist, dass bei
PTSD-Betroffenen keines der oben genannten Leitsymptome
manifest sein muss, da viel bereits über langjährige
Symptomkompetenz verfügen bzw. durch Therapieerfahrung
Kontrolle über diese erzielt haben, ein übermäßig
bemüht-rücksichtsvoller also seitens des Klienten auch als
positiv diskriminierend bzw. stigmatisierend erlebt werden
kann!}
{}
{}
{}
{}

\Text{Andauernde Persönlichkeitsänderung nach Extrembelastung (F62.0)}
{Eine andauernde, wenigstens über zwei Jahre bestehende
Persönlichkeitsänderung kann einer Belastung katastrophalen
Ausmaßes folgen. Die Belastung muss extrem sein, dass die
Vulnerabilität der betreffenden Person als Erklärung für die
tief greifende Auswirkung auf die Persönlichkeit nicht in
Erwägung gezogen werden muss. Die Störung ist durch eine
feindliche oder misstrauische Haltung gegenüber der Welt,
durch sozialen Rückzug, Gefühle der Leere oder
Hoffnungslosigkeit, ein chronisches Gefühl der Anspannung
wie bei ständigem Bedrohtsein und Entfremdungsgefühl,
gekennzeichnet. Eine posttraumatische Belastungsstörung
(F43.1) kann dieser Form der Persönlichkeitsänderung
vorausgegangen sein.

Mögliche Auslöser können sein:
\begin{itemize}
  \item andauerndes Ausgesetztsein lebensbedrohlicher
  Situationen z.\,B. durch wiederholte Gewalterfahrungen,
  Missbrauch etc.
  \item andauernder Gefangenschaft mit unmittelbarer Todesgefahr
  \item Folter
  \item Katastrophen
\end{itemize}
\subparagraph{Leitsymptome}
\begin{itemize}
  \item erhöhte Anspannung
  \item Deutlich Veränderte kognitive
  Schemata/Kontrollüberzeugungen, z.\,B. deutliches
  Misstrauen gegenüber Mitmenschen, starkes Gefühl  der
  Entfremdung und ständiges sich Bedrohtfühlen, evtl. auch
  resignative bzw. zynische Grundhaltung.
  \item Diagnosestellung erst nach zwei Jahren (ab dem
  Zeitpunkt der Traumatisierung) bei anhaltender
  Persönlichkeitsänderung möglich.
\end{itemize}

\subparagraph{Folgerungen für den Sanitätsdienst} Sozial auffälliges
Verhalten als Folge des prinzipiellen Misstrauens bzw.
Bedrohtfühlens können den Umgang in der Versorgungs-
Betreuungssituation erschweren. Es mag hilfreich erscheinen,
dieses Verhalten nicht auf sich zu beziehen/persönlich zu
nehmen, sondern als Ausdruck der zugrundeliegenden
Persönlichkeitsänderung zu deuten.}
{}
{}
{}
{}


\Text{Retraumatisierung}
{Darunter ist eine \EE{mögliche Reaktivierung} einer bereits
\EE{latenten oder abgemilderten PTSD-Symptomatik} durch die
\EE{Konfrontation mit traumaassoziierten Reizen} und oder
erneute Traumatisierungen gemeint, was eine erneute
Manifestation von Symptomen bzw. einen Symptomanstieg zur
Folge haben kann und jedenfalls von Betroffenen als
belastend erlebt wird. Symptome einer Retraumatisierung
ähneln zum Teil auch jenen einer akuten Belastungsreaktion,
da die Traumasituation subjektiv wiedererlebt wird.
Verhaltensregeln, um Retraumatisierungen zu vermeiden, siehe
\enquote{Posttraumatische Belastungsstörung, Folgerungen für
den Sanitätsdienst}.}
{}
{}
{}
{}

Bezeichnet das Phänomen, dass von Traumafolgestörungen betroffene Menschen durch die Konfrontation mit traumaassoziierten Reizen (Sinneswahrnehmungen, Gegenstände, Orte, Erinnerungen, etc.) erneut traumatisiert werden können, d.h. eine bereits abgemilderte Symptomatik wieder aufkeimen/sich intensivieren kann, was von Betroffenen als sehr unangenehm und belastend erlebt wird. Im Sinne des Schutzes Betroffener lässt sich daraus eine erhöhte Vorsicht ableiten (siehe Posttraumatische Belastungsreaktion, Folgerungen für den Sanitätsdienst, v.a. Rücksichtnahme auf Vermeidungsverhalten).


\Text{Sekundärtraumatisierung}
{Unter sekundärer Traumatisierung - auch induzierte oder
stellvertretende Traumatisierung (vicarious traumatisation)
genannt -  bezeichnet eine \EE{Veränderung des inneren
Erlebens als Folge mitfühlenden Engagements mit Traumaopfern
und/oder traumatischen Inhalten}. Für Sanitäter könnte dies
durch den häufigen Umgang mit traumatisierten Menschen
tragend werden, nicht nur durch die tatsächliche
Konfrontation mit der Extremsituationen, sondern auch durch
die in der Betreuungssituation stattfindenden Gespräche. Die
Fähigkeit sich in seiner Rolle als professioneller Helfer
auch ausreichend abgrenzen und von erlebten
Belastungssituation Betroffener/erzählten Traumainhalten
distanzieren zu können, stellt einen wichtigen Schutz dar
(siehe \enquote{Exkurs: Rollenkonflikt} bzw. Abschnitt
\enquote{Ressourcen}).

\Text{Sekundärtraumatisierung}
{Unter sekundärer Traumatisierung - auch induzierte oder
Das Phänomen der stellvertretende Traumatisierung (vicarious traumatisation) beschreibt Veränderungen des inneren
Erlebens von helfenden Professionisten als Folge ihres mitfühlenden Engagements bzw. des häufigen Umgangs mit Traumaopfern sowie 
mit Menschen in Extremsituationen. Dies umfasst auch die in der Betreuungssituation stattfindenden Gespräche. Die
Fähigkeit, sich die Rolle als professioneller Helfer bewusst zu machen und sich ausreichend von erlebten Belastungssituation Betroffener/erzählten Traumainhalten distanzieren zu können, stellt eine wichtige und schützende Kompetenz dar. Ein einfühlender Umgang ist nicht als grenzenlose Gefühlsübernahme und -ansteckung zu verstehen! 


Mögliche Anzeichen einer Sekundärtraumatisierung sind
\begin{itemize}
  \item Empathiemüdigkeit, „Seelenverlust" der Professionellen
  \item Psychische Ansteckung, Gefühle der Verstörung,
  Konfrontation mit eigenen Anteilen von Grausamkeit,
  Machtmissbrauch und Destruktivität.
  \item Zusammenbruch der
  persönlichen Bewältigungskompetenzen bzw. des
  Kohärenzerlebens
  \item Verunsicherung der eigenen Identität,
  Bedrohung der Bedürfnisse nach Sicherheit und Geborgenheit
  in der Welt
  \item Verlust von Lebensfreude und Arbeitsfreude
  \item Konfrontation mit der eigenen Verletzlichkeit und der
  Realität traumatischer Verhältnisse
  \item Desillusionierung,
  Hoffnungslosigkeit, Ohnmacht, Trauer
  \item Existentielle
  Krise und Erschütterung, Zweifel am Sinn von Leben, Leiden,
  Tod
\end{itemize}}
{}
{}
{}
{}

\begin{itemize}
  \item  Konfrontation mit der eigenen Verletzlichkeit und der
  Realität traumatischer Verhältnisse, Konfrontation mit eigenen Anteilen von Grausamkeit, Machtmissbrauch und Destruktivität.
  \item Verlust von Lebensfreude und Arbeitsfreude, Empathiemüdigkeit, „Seelenverlust“, „Psychische Ansteckung“, Gefühle der Verstörung,
  \item Bedrohung der Bedürfnisse nach Sicherheit und Geborgenheit
  in der Welt, Verunsicherung der eigenen Identität, Zusammenbruch der
  persönlichen Bewältigungskompetenzen bzw. des
  Kohärenzerlebens
  \item Desillusionierung, Hoffnungslosigkeit, Ohnmacht, Trauer, Existentielle Krise und Erschütterung, Zweifel am Sinn von Leben, Leiden, Tod
\end{itemize}}


\A{Keywords: Psychische/physische Ebene,  erhöhte
Vulnerabilität, Frustrationserleben (Reizbarkeit, gedrückte
Stimmung), Zynismus, pathologische Phänomene v.a. Burnout
und Traumafolgen (akute Belastungsreaktion, PTSD, Andauernde
Persönlichkeitsänderung nach Extrembelastung),
Retraumatisierung, Sekundärtraumatisierung.}

\section{Stressbewältigung}

\subsection{Das Copingkonzept}

\Definition{Stressbewältigung (\T{Coping}) ist der \Speech{\enquote{Prozess der
Handhabung jener externen oder internen Anforderungen, die vom Individuum als
die eigenen Ressourcen beanspruchend oder übersteigend
bewertet werden}}.\ZFootnote{Lazarus \&
Folkman (1984)} 
Bei \T{Copingstrategien} handelt es sich um \Speech{\enquote*{alle
\EE{kognitiven, emotionalen und behaviouralen Anstrengungen}, die dazu
dienen \EE{Belastung und Stress zu
bewältigen}}}.\ZFootnote{Trautmann-Sponsel (1988)}}


Nach Lazarus hat Coping hauptsächlich folgende
\EE{Aufgaben}:
\begin{itemize}
  \item Den Einfluss schädigender Umweltbedingungen zu
  reduzieren und die Aussicht auf Erholung zu verbessern.
  \item Negative Ereignisse oder Umstände ertragbar machen
  bzw. den Organismus an sie anpassen.
  \item Ein positives Selbstbild aufrechterhalten.
  \item Das emotionale Gleichgewicht sichern.
  \item Befriedigende Beziehungen zu anderen Personen fortsetzen.
\end{itemize}

Welche Coping-Reaktionen jemand in einer bestimmten
Situation wählt, hängt von diversen Faktoren ab:
\begin{itemize}
  \item allgemeiner Gesundheitszustand
  \item Grad der psychischen und/oder physischen Belastung
  \item Bereich, von dem die Anforderungen ausgehen
  \item Zeitfaktor
  \item Frühere Erfolge bzw. Misserfolge bei ähnlich
  strukturierten Anforderungssituationen
  \item Grad der subjektiven Bedeutsamkeit
  \item Persönliche Präferenzen aufgrund von Sozialisationserfahrungen
\end{itemize}


\Text{Aspekte des Copings}
{\Z{Coping stellt ein \EE{prozesshaft-dynamisches Geschehen}
dar, mit außerordentlich vielen Variationsmöglichkeiten.
Diverse Verhaltensweise, die als Reaktion auf Stress gesetzt
werden, können als Coping-Reaktion interpretiert werden,
weshalb es schwierig ist einen systematischen Überblick zu
geben.

Grundsätzlich können Bewältigungsversuche belastenden
Ereignissen entweder \EE{ereignisbezogen} (z.\,B.
\Speech{\enquote{Für dieses Problem muss es doch eine Lösung
geben}}) oder \EE{selbstzentriert} (\Speech{\enquote{Hätte ich doch nur
nicht so gehandelt \ldots}}).

Um eine Stresssituation zu meistern, suchen manche nach
\EE{sozialer Unterstützung}, holen sich Hilfe bzw.
Ratschläge bei Familie,  Freunden oder Kollegen, während
andere aus Angst vor Ansehens- u. Prestigeverlust versuchen
die Situation \EE{alleine} zu bewältigen.

Coping kann rein \EE{innerpsychisch} ablaufen, d.\,h. es
bleibt auf Gefühle und Gedanken begrenzt, z.\,B. wenn
versucht wird, den zunächst als hoch eingeschätzten
Bedeutungsgehalt einer Situation oder deren Folgen in einer
Art Uminterpretation herunterzuspielen, sich von der
Anforderung gedanklich zu distanzieren bzw. sich selbst aus
der Verantwortung zu nehmen.

Zuguterletzt kann es sich in direkten Handlungen
manifestieren: hier reicht die Spanne möglicher Reaktionen
von \EE{Flucht und Vermeidung bis hin zu einer
selbstbewussten und selbstbestimmten aktiven
Auseinandersetzung} mit der Herausforderung: Man wägt
zunächst seine Handlungsmöglichkeiten ab, das Problem mit
den Ressourcen, die einem selbst zur Verfügung stehen, zu
lösen. Erweisen sich diese als ungenügend, kann versucht
werden, den eigenen Informationsstand zu auszuweiten oder
aber nach Möglichkeiten gesucht werden, die eigenen
Fähigkeiten zu erweitern.}}
{}
{}
{}
{}


\Text{Arten von Coping}
{Folgende Arten von Coping lassen sich unterscheiden\cite{ZimbardoPsych6}:

\begin{itemize}
  \item \EE{Problemorientiertes Coping}
  Dieses umfasst alle Strategien des direkten Umgangs mit dem
  Stressor, d.\,h. hier versucht der Betroffene durch
  \EE{Informationssuche, direktes Handeln oder auch
  Unterlassen von Handlungen} Belastungssituationen zu
  bewältigen oder sich den Gegebenheiten anzupassen. Diese
  Copingstrategie orientiert sich an der Ebene der Situation
  bzw. des Reizes, z.B. Großschadensereignis als massiver Stressor → Informationssuche („Was ist passiert, wer ist wie betroffen?“) um Klarheit zu schaffen und weitere Schritte planen zu können, gezielte Behandlungsinterventionen nach Plan abarbeiten, Hektische und komplizierte Handlungen vermeiden.
  \item \EE{Emotionsorientiertes Coping}
  Hierbei handelt es sich um innerpsychisches Coping, bei dem
  hauptsächlich versucht wird, die durch die Situation
  hervorgerufenen emotionalen Spannungen und unangenehmen
  Gedanken zu verändern (\EE{=Emotionsregulation}). Z.B. durch Umdeutung von Gefühlen (Stressgefühl und Ängstlichkeit sind natürliche Reaktionen und machen mich handlungsfähig), positive Selbstinstruktion (Ich schaffe das, ich habe ähnliche Situationen auch schon gemeistert)
  \item \EE{Bewertungsorientiertes Coping}
  Dieses setzt in der Phase der Bewertung bzw. Neubewertung
  nach Lazarus´Stresstheorie an: einerseits da, wo es um die
  \EE{Interpretation einer Situation als Bedrohung, Belastung
  oder eben Herausforderung} geht. Zum anderen in der Phase
  der Neubewertung, indem die betroffene Person ihr
  \EE{Verhältnis zur Umwelt kognitiv neu bewerten kann}, um so
  einen adäquaten Umgang zu finden. Hauptziel ist somit eine
  \EE{Belastung eher als Herausforderung} zu sehen, weil damit
  der jeweiligen Lebensumstand positiv bewertet wird und
  dadurch eher Handlungsmöglichkeiten bzw.
  Bewältigungsressourcen frei werden, um angemessen zu
  reagieren. Dies erfordert im Weitern das Auffinden konkreter
  Problemlösungsansätze (siehe problemorientiertes Coping) und
  somit die Kombination unterschiedlicher
  Bewältigungsstrategien. Z.B.: Neubewertung/-einschätzung nach Abklingen erster alarmierender, lähmender Emotionen: „Die Situation scheint ernst, aber nicht aussichtslos.“ „Ich kann nicht alle ungünstigen Bedingungen beeinflussen, aber x,y und z könnte ich verändern...“
\end{itemize}

}
{}
{}
{}
{}




\subsection{Die Theorie der Ressourcenerhaltung  }

Nach der Theorie der Ressourcenerhaltung \EE{tritt Stress
dann auf, wenn Ressourcen bedroht werden, verloren gegangen sind oder
fehlinvestiert wurden}, d.\,h. er sieht nicht Herausforderung,
Bedrohung oder Schädigung, sondern das Erleben von Verlust
als zentrale Eigenschaft von
Stress.\ZFootnote{Die Theorie von Hobfoll (zitiert nach
\cite{StarkeMedSozAspProbBew00}) stellt eine modernere
Alternative zu der Theorie von Lazarus dar bzw. kann auch
als deren Komplement angesehen werden. Sie stellt
nicht die Einschätzungen/Bewertungsprozesse, sondern die
Stressbewältigung und die Beweggründe dafür in den Fokus,
d.\,h. es geht Hobfoll hauptsächlich um Motivation und
Coping, wofür er den Begriff der Ressourcen prägt (siehe
unten).

Er nimmt an, dass \EE{Individuen danach streben, was
sie wertschätzen}, es erhalten und Verluste vermeiden bzw.
minimieren wollen, wobei für ihn die Vermeidung von
Verlusten als stärkeres Motiv als das Streben nach
Zugewinnen gilt.

\EE{Ableitungen aus seinen grundlegenden Annahmen}:
\begin{inparaenum}
  \item Ressourcen sind nicht zufällig verteilt. Menschen
  spielen eine \EE{aktive Rolle} darin, \EE{über
  welche Ressourcen sie verfügen und wie geschickt sie diese
  zum Einsatz bringen}.
  \item Menschen \EE{setzen Ressourcen} dazu ein, \EE{um
  andere Ressourcen zu erhalten}. Wer also über viele und gute
  Ressourcen verfügt, ist gegenüber den Schicksalsschlägen des
  Lebens nicht nur besser gewappnet, sondern kann auch immer
  mehr investieren, um immer mehr zu gewinnen.
  \item \EE{Coping verursacht Kosten}, da Ressourcen zur
  Sicherung anderer Ressourcen gebraucht werden.
  \item \EE{Gewinn und Verlust} von Ressourcen \EE{erzeugen
  positive und negative Spiralen}.
\end{inparaenum}

Die Theorie postuliert eine Wechselwirkung von Ressourcen,
Bedürfnissen, Belastungen, Werten, Zeit und Wahrnehmungen.
Daraus folgt eine mehr oder weniger gelungene Passung
zwischen den jeweiligen Anforderungen der Umwelt und den
individuellen Reaktionen. Wie Stress bewältigt wird, hängt
demnach von dem Grad dieser Passung innerhalb eines
bestimmten ökologischen Kontextes ab.}

% \EE{Klassifikation von Ressourcen}

\Text{Ressourcen}
{Ressourcen sind \EE{Objekte, persönliche Charakteristika,
Bedingungen und Energien, die vom Individuum wertgeschätzt
und als unterstützend erlebt} werden. Sie stellen
vereinfacht dargestellt das Gegenkonzept zu Stressoren dar,
als Faktoren, die das persönliche Kompetenzerleben,
Sicherheitsgefühl und Wohlbefinden erhalten, schützen oder
sogar steigern. Sie gelten daher auch als Schutzfaktoren.


Ressourcen können sein \ldots
\begin{itemize}
  \item \EE{Gegenstände}, z.\,B. Kleidung, Wertsachen, Autos,
  Nahrungsmittel, und andere Objekte dieser Art.
  \item \EE{Bedingungen}, d.\,h. nichtmaterielle Ressourcen,
  z.\,B. Qualifikationen, Ehe, Karriere, Staatsbürgerschaft,
  Beamtenstatus, Autonomie, Beteiligung an
  Entscheidungsprozessen etc.
  \item \EE{Persönlichkeitsmerkmale}: stabile Fähigkeiten,
  Fertigkeiten, persönliche Überzeugungen, z.\,B. Intelligenz,
  Selbstwirksamkeit, Empathie, soziale Verantwortung.
  Optimismus, aber auch erlebte Gefühle wie Erfahrung von
  Liebe, Zuneigung und Wertschätzung durch andere etc.
  \item \EE{Energien}, darunter verstanden werden nicht von
  sich aus wertvolle Ressourcen, sondern solche, die Zugang zu
  anderen Ressourcen ermöglichen, d.\,h. sie sind konvertierbar
  in Gegenstände, Bedingungen oder andere Energien, die man
  gerade in bestimmten Stresssituationen benötigt, z.\,B.
  Geld, Zeit, Wissen
  ...die als wertvoll und nützlich erscheinen.
\end{itemize}}
{}
{}
{}
{}


\EE{Über je mehr Ressourcen man verfügt, desto  mehr
Belastungen können bewältigt werden, ABER: dies macht einen
nicht unverwundbar! Es wird immer Situationen, die so
intensiv und außergewöhnlich erlebt werden, dass sie die
persönlich zur Verfügung stehenden Bewältigungsmittel eines
Menschen überfordern können und sei er noch so reich an
Ressourcen/Schutzfaktoren.}

\subsection{Das Konzept des Kohärenzgefühls}

\Z{Antonowsky (1997) definiert das Kohärenzgefühl als eine
\EE{globale Orientierung}, die ausdrückt, in welchem Ausmaß
man ein durchdringendes, andauerndes und dennoch dynamisches
\EE{Gefühl des Vertrauens} hat, dass
\begin{itemize}
  \item die \EE{Stimuli}, die sich im Verlauf des Lebens aus
  der inneren und äußeren Umgebung ergeben, \EE{strukturiert,
  vorhersehbar und erklärbar} sind.
  \item einem die \EE{Ressourcen zur Verfügung} stehen werden,
  um den \EE{Anforderungen}, die diese Stimuli stellen, \EE{zu
  begegnen}
  \item diese Anforderungen \EE{Herausforderungen} sind, die
  \EE{Anstrengung und Engagement lohnen}.
\end{itemize}
In anderen Worten handelt: es sich als um eine
grundoptimistische Haltung, Anforderungen erkennen und
bewältigen zu können, einem Grundvertrauen in die eigenen
Kompetenzen sowie die Überzeugung, dass die
Herausforderungen anzunehmen, bedeutsam und sinnvoll
erscheint. Dies lässt sich in die drei Komponenten
\EE{Verstehbarkeit (1. Punkt), Handhabbarkeit (2. Punkt) und
Bedeutsamkeit (3. Punkt)} zusammenfassen.

Je nach Ausprägung in den o.\,g. Anteilen des
Kohärenzgefühls lassen sich div. Konfigurationen vermuten.
Im Folgenden sei
ausschließlich von einem starken -- d.\,h. auf allen drei
Ebenen hoch ausgeprägten -- und einem schwachen -- d.\,h. in
allen drei Ebenen niedrig ausgeprägten -- Kohärenzgefühl
eingegangen.

\EE{Ableitungen} in Bezug auf die Wirkungsweise des
Kohärenzgefühls:
\begin{itemize}
  \item \EE{Auswahl gesundheitsfördenden Verhaltens}:
  Personen mit hohem Kohärenzgefühl interpretieren weniger
  Stimuli als stresshaft und weichen Stressoren, die schwer
  handzuhaben sind, eher aus.
  \item \EE{Erfolgreicherer Umgang mit Stressoren}: Personen
  mit hohem Kohärenzgefühl \enquote{copen} erfolgreicher, was zu
  Spannungsreduktion und Vermeidung von Schädigung führe.
\end{itemize}
Inwieweit sich diese von Antonovsky postulierten
Grundhaltungen entwickeln, modifizieren oder verändern
lassen ist nicht eindeutig ergründet, zu schreibt der Autor
dem Konstrukt stark dispositionellen Charakter zu.
Augenscheinlich ergeben sich \EE{für den Sanitätsdienst
nachstehende Folgerungen} ableiten:

\begin{itemize}
  \item Verstehbarkeit kann durch Informations- und
  Wissensvermittlung, durch Fortbildung gefördert werden. Je
  ausgeprägter das Wissen über mögliche Betreuungs- und
  Versorgungssituationen, Unfallfolgen, Krankheitsbilder und
  Zustandsbilder ist, desto eher wird sich die
  Informationsfülle reale Notfallsituationen strukturieren,
  priorisieren und auf das Wesentliche reduzieren lassen -
  also verstehbarer.
  
  \item Ebenso lässt sich im Bezug auf die Handhabbarkeit
  argumentieren: Kompetenzen durch praktisches Üben,
  Trainings, Schulungen etc. zu erweitern, kann das
  persönliche Kompetenzerleben steigern und schwierige
  Situationen als bewältigbar erleben.
  
  \item Hinsichtlich des Aspekts der Bedeutsamkeit ist
  anzumerken, dass diese v.a. unter div. o.b.
  Stressauswirkungen und insbesondere bei pathologischen
  Entwicklungen (Burnout, Traumafolgen) und oder
  Sekundärtraumatisierung sehr stark in Mitleidenschaft
  gezogen werden kann.
\end{itemize}
}
\A{Keywords: Stressbewältigung,  Copingstrategien (Exkurs:
Proaktives Coping), Theorie der Ressourcenerhaltung,
Kohärenzgefühl, (Resilienz), (Abwehrmechanismen).}

\subsection{Das Konzept des Kohärenzgefühls}

Unter dem Begriff Kohärenzgefühl versteht Antonovsky eine
\EE{globale Orientierung}, die ausdrückt, in welchem Ausmaß
man ein durchdringendes, andauerndes und dennoch dynamisches
\EE{Gefühl des Vertrauens} hat, dass
\begin{itemize}
  \item die \EE{Eindrücke/Erlebnisse}, die sich im Verlauf des Lebens aus
  der inneren und äußeren Umgebung ergeben, \EE{eine Ordnung haben,
  vorhersehbar und erklärbar} sind (=Verstehbarkeit).
  \item einem persönlich die \EE{Ressourcen zur Verfügung} stehen werden,
  um den \EE{Anforderungen}, die diese Eindrücke/Erlebnisse stellen, \EE{zu begegnen} (=Handhabbarkeit)
  \item diese Anforderungen \EE{Herausforderungen} sind, die
  \EE{Anstrengung und Engagement lohnen} (=Bedeutsamkeit).
\end{itemize}
In anderen Worten handelt: es sich als um eine
generelle grundoptimistische Haltung, einem Grundvertrauen Anforderungen erkennen und mithilfe persönlicher Kompetenzen bewältigen zu können sowie die Überzeugung, dass die Herausforderungen anzunehmen, sinnvoll und bedeutsam erscheint. 

Je nach Ausprägung in den o.\,g. Anteilen des
Kohärenzgefühls lassen sich div. Konfigurationen vermuten.
Im Folgenden sei ausschließlich von einem starken -- d.\,h. auf allen drei Ebenen hoch ausgeprägten -- und einem schwachen -- d.\,h. in
allen drei Ebenen niedrig ausgeprägten -- Kohärenzgefühl
eingegangen.

\EE{Ableitungen} in Bezug auf die Wirkungsweise des
Kohärenzgefühls:
\begin{itemize}
  \item \EE{Auswahl gesundheitsfördenden Verhaltens}:
  Personen mit hohem Kohärenzgefühl interpretieren weniger
  Eindrücke als stresshaft (z.B. weil sie sie als bewältigbar einschätzen) und weichen Stressoren, die schwer
  handzuhaben sind, eher aus.
  \item \EE{Erfolgreicherer Umgang mit Stressoren}: Personen
  mit hohem Kohärenzgefühl \enquote{verarbeiten} Stressoren erfolgreicher, was zu
  Spannungsreduktion und Vermeidung von Schädigung führe.
\end{itemize}
Inwieweit sich diese von Antonovsky postulierten
Grundhaltungen entwickeln, modifizieren oder verändern
lassen ist nicht eindeutig ergründet.

Augenscheinlich lassen sich \EE{für den Sanitätsdienst
nachstehende Folgerungen} ableiten:

\begin{itemize}
  \item Verstehbarkeit kann durch Informations- und
  Wissensvermittlung, z.B. durch Fortbildungen, Schulungen,
  Selbststudium etc. gefördert werden. Je ausgeprägter das
  Wissen über mögliche Betreuungs- und
  Versorgungssituationen, Unfallfolgen, Krankheits- und
  Zustandsbilder ist, desto eher wird sich die
  Informationsfülle und -komplexität realer Notfallsituationen strukturieren,
  priorisieren und auf das Wesentliche reduzieren lassen --
  also verstehbarer.
  
  \item Ebenso lässt sich im Bezug auf die Handhabbarkeit
  argumentieren: Kompetenzen durch praktisches Üben,
  Trainings, Schulungen etc. zu erweitern, kann das
  persönliche Kompetenzerleben steigern, wodurch herausfordernde
  Situationen eher als bewältigbar erlebt werden.
  
  \item Hinsichtlich des Aspekts der Bedeutsamkeit ist
  anzumerken, dass dieser v.\,a. unter divierse o.\,g.
  Stressauswirkungen und insbesondere bei pathologischen
  Entwicklungen (Burnout, Traumafolgen) und
  Sekundärtraumatisierung sehr stark in Mitleidenschaft
  gezogen werden kann. Sowohl sehr niedrige (kein
  Sinnerleben mehr in der Tätigkeit) als auch hohe
  Bedeutsamkeit (Übereifer, Überidentifikation,
  Überverantwortlichkeit)kann Stressfolge als auch -wirkung
  sein und könnte als persönlicher „Thermometer“ einen
  Hinweis auf Stressfolgen (siehe auch:
  Sekundärtraumatisierung, Burnout) darstellen.
\end{itemize}

\A{Keywords: Stressbewältigung,  Copingstrategien (Exkurs:
Proaktives Coping), Theorie der Ressourcenerhaltung,
Kohärenzgefühl, (Resilienz), (Abwehrmechanismen).}

\section{Grundsätze der Stressbewältigung im Sanitätsdienst, Maßnahmen zur Verhütung und Verminderung von Beanspruchungsfolgen }


Im Sanitätsberuf lassen sich nach 
bestimmte Situationen identifizieren, die von vielen Professionisten als besonders unangenehm
bewerten\ZFootnote{Ungerer
(1999), zitiert nach \cite{Karutz200912}}, z.\,B. Einsätze,
\begin{itemize}
  \item bei denen \EE{Kinder} schwer verletzt worden sind,
  \item bei denen die \EE{Helfer selbst} in akute Gefahr
  geraten,
  \item bei denen \EE{Kollegen verletzt oder getötet}
  werden,
  \item bei denen die Opfer den Helfern \EE{persönlich
  bekannt} sind,
  \item bei denen die Helfer einen anderen \EE{persönlichen
  Bezug} zu den Opfern herstellen,
  \item in denen Helfer mit besonders schrecklichen
  Anblicken konfrontiert werden,
  \item die sehr lang andauern und auch körperlich besonders
  beanspruchend sind,
  \item die besonders tragisch sind, sowie
  \item um Einsätze, die besonders dramatisch bzw.
  \EE{dynamisch} und unter besonders großem \EE{Zeitdruck}
  ablaufen.
\end{itemize}

Häufig werden diese Situationen von einer psychophysischen
Übererregung -- im Sinne des oben beschriebenen
Erregungsniveaus (siehe \enquote{Aktivierung}) -- einhergehen,
sodass das jeweilige Notfallgeschehen und die durchzuführenden
Maßnahmen nicht mehr als durchaus positiv bewertete
Herausforderung, sondern eher negativ, d.\,h. als eine
Belastung bzw. Überforderung erlebt werden.


\Text{Indikation der psychologischen Selbsthilfe}
{
Verschiedene Symptome weisen auf eine hohe Belastung von
professionellen Helfern hin (siehe Tabellen) und können
Hinweise bzw. Warnsignale darstellen, entsprechend zu
intervenieren (siehe Psychologische Selbsthilfe). Um diese
möglichst frühzeitig zu erkennen sollten
\Speech{\enquote{Einsatzkräfte [...] daher auch unmittelbar im Einsatzgeschehen auf ihre
eigene Befindlichkeit achten und ihre psychische Situation
angemessen beurteilen, ohne sich selbst zu über- oder zu
unterschätzen }}\ZFootnote{Lasogga \& Karutz 2005, zitiert nach
Karutz 2009}.}
{}
{}
{}
{}





\subsection{Präventive Strategien}


Die folgenden beschriebenen Strategien sollen  vorwiegend
darauf abzielen, \Speech{\enquote{die Handlungsfähigkeit eines
Helfers in einer absoluten Ausnahmesituation und nur für
eine relativ kurze Zeit aufrecht zu erhalten}} und dienen nicht generelle als langfristige Methoden zur
Stressbewältigung \cite{Karutz200912}.


\Text{Einstellung}
{Eine positive Grundeinstellung an Einsätze
kann helfen, die Hoffnung auf Erfolg zu stützen,
Erwartungsängste zu mildern und die Situation als
Herausforderung anzusehen. Sollten Einsatzmeldung
beunruhigend wirken\ZFootnote{Ungerer, 1999 zitiert nach Karutz,
2009}, kann das Bewusstmachen möglichst angenehmer
Vorerfahrungen z.\,B. Erinnerungen an ähnliche Einsätze, die
bereits erfolgreich bewältigt werden konnten nach dem Motto
\SpeechTextMargin{\enquote{Bisher haben wir das stets gut bewältigt, also
werden wir es auch diesmal hinbekommen!}} In anderen Fällen
kann zumindest der Gedanke \SpeechTextMargin{\enquote{Irgendwie werden wir auf
jeden Fall helfen können}} helfen.\ZFootnote{Lasogga \& Karutz 2005,
zitiert nach \cite{Karutz200912}}

\II{Erwartungshaltung}: Überhöhte Erwartungen
an das eigene Handeln können Anspannung und zusätzlichen
Druck erzeugen. Extrem übersteigerte Anspruchshaltungen
(\SpeechTextMargin{\enquote{Ich will immer alles perfekt machen!}}) können
somit die Fehlerquote erhöhen (siehe Yerkes-Dodson-Gesetz). Zu
akzeptieren, dass kleinere Unzulänglichkeiten passieren
können, kann Druck nehmen und ermöglicht auch ein
\II{Fehlerbewusstsein}: nämlich Fehlerquellen zu erkennen,
zu analysieren und zu optimieren. \SpeechTextMargin{\enquote{Ohnehin verläuft
kein Einsatz perfekt, und aus dem, was suboptimal gewesen ist,
kann und sollte man eben für die Zukunft lernen}}\ZFootnote{Karutz et al., 2003 zitiert nach \cite{Karutz200912}}.

Es wird immer Notfallsituationen geben, die trotz hoch engagierter und
kompetenter Hilfeleistung kaum oder nicht beeinflussbar
sind. Weiters kann die Entstehung unbegründeter
Schuldgefühle und Selbstvorwürfe bzw. einem Gefühl der
Überverantwortlichkeit durch eine realistische Einschätzung
der Handlungsmöglichkeiten vorgebeugt werden.
\Speech{\enquote{Rettungsdienstmitarbeiter können das tun, was
menschlich bzw. fachlich unter den gegebenen Bedingungen möglich ist –
sie sollten sich aber auch zugestehen, dass dies in einigen
Fällen nun einmal nicht ausreicht: \EE{70 bis 95 Prozent der
Reanimationsversuche enden beispielsweise erfolglos}, so
sehr man sich dies sicherlich anders wünschen würde}}\ZFootnote{Baskett,
Stehen \& Bossaert, 2006, zitiert nach \cite{Karutz200912}}.}
{}
{}
{}
{}

\Text{Bewertung von Stress und Angst}
{Ein gewisses Maß an
Anspannungserleben ist -- wie in vorangegangenen Kapiteln
beschrieben -- erforderlich, ja sogar unumgänglich, um eine
grundsätzliche Bereitschaft zur Leistungsfähigkeit
herzustellen. Im Einsatz einen gewissen Grad an Aufregung
wahrzunehmen sollte somit nicht gleich einen Grund zur Sorge
oder gar Scham darstellen. Aktivierung bewirkt über die
Ausschüttung von Stresshormonen entsprechend des
Yerkes-Dodson-Gesetzes zunächst eine Leistungssteigerung im
Organismus: \Speech{\enquote{Die Konzentrationsfähigkeit steigt,
die Muskulatur wird besser durchblutet usw. Es werden also die
Grundvoraussetzungen geschaffen, auf die Anforderungen der
Versorgungssituation adäquat reagieren zu können}}, d.\,h.
\SpeechTextMargin{\enquote{{\ldots} eine gewisse Anspannung ist also
keineswegs negativ zu bewerten, sondern – ganz im Gegenteil – für die persönliche
Handlungsfähigkeit hilfreich und notwendig}}\ZFootnote{Lasogga \&
Karutz 2005, zitiert nach \cite{Karutz200912}}.  Ebenso
stellt das Erleben einer gewissen Angst ein evolutionär verankertes
Frühwarnsystem dar, das Wachsamkeit und Vorsicht fördert und
daher nicht generell als negativ bewertet werden sollte,
sofern sie nicht lähmend wirkt.}
{}
{}
{}
{}





\subsubsection{Kurzfristige Strategien}

\Text{Aufgaben in Gedanken durchgehen}
{Die einzelnen
Maßnahmen zu strukturieren, nach Prioritäten zu ordnen und
die wichtigsten Handlungsrichtlinien (s. unten)
durchzudenken kann das eigene Sicherheitsgefühl stärken und
Halt geben:

\begin{itemize}
  \item \Speech{\enquote{Eigenschutz geht immer vor!}}
  \item \Speech{\enquote{Tue nur das, was Du sicher beherrschst}}
  \item \Speech{\enquote{Behandle zuerst, was zuerst tötet}}
  \item \Speech{\enquote{Erst einmal Platz schaffen, damit man gut
  arbeiten kann!}} usw.\ZFootnote{Karutz , 2003 zitiert nach
  \cite{Karutz200912}}.
\end{itemize}

Ebenso kann das in  Erinnerung rufen wichtiger
Behandlungsmodule z.\,B. das ABC der Reanimation, die
einzelnen Phasen der Traumaversorgung oder die 4-S-Sätze zur
Psychischen Ersten Hilfe. \Speech{\enquote{Ebenso zweckmäßig ist
das Abarbeiten von Checklisten oder Algorithmen: Je mehr diese
in der Ausbildung verinnerlicht worden sind, umso eher
werden sie auch in einer Stresssituation noch angewendet
bzw. automatisiert abgerufen werden können}}
(\cite{Karutz200912}).}
{}
{}
{}
{}



\Text{Mentale Vorbereitung, Alternativen planen (\enquote{Plan B})}
{Gemeint ist damit, alternative Handlungsmöglichkeiten zu
überlegen bzw. vorauszuplanen, eine Art Spickzettel mit
möglichen Handlungen, die sich auch dann noch ausführen
lassen, wenn alle andere Interventionen bereits ausgeschöpft
wurden, um das Gefühl der Handhabbarkeit zu fördern. Ein
solcher persönlicher Notfallplan kann laut \Study{Karutz}
\cite{Karutz200912} vor möglichen Gefühlen der Hilflosigkeit und Ohnmacht schützen,
er führt folgende Beispiele an:
\begin{itemize}
  \item sich einige Meter vom unmittelbaren Ort des Geschehens entfernen, beispielsweise um einen benötigten Ausrüstungsgegenstand aus dem Einsatzfahrzeug zu holen - Räumliche Distanz ermöglicht auch psychische Distanz.
  \item Kollegen um Unterstützung zu bitten.
  \item In Großschadenslagen eventuell kleinere Pausen
  einlegen, um etwas zu trinken oder zu essen.
  \item Sich notfalls ablösen zu lassen.  Dieses sollte nicht
  \Speech{\enquote{als persönliches Versagen interpretiert werden,
  eher handelt es sich um den Ausdruck realistischer
  Selbsteinschätzung}}\ZFootnote{Karutz et al., 2003 zitiert nach
  Karutz, 2009}.
\end{itemize}}
{}
{}
{}
{}




\Text{Kognitive Umbewertung}
{Wie aus dem transaktionalen
Stressmodell ableiten lässt stellen Bewertungsprozesse bei
der Verarbeitung von Stress einen wesentlichen Faktor dar:
dies lässt sich auch auf Notfallsituationen anwenden, im
Sinn eines positiven Umdeutens \cite{BengelPsyNotfallmedRd2-PsyBelRdPers}  durch
problemschwächende Interpretationen (z.\,B. \enquote{Eine solche
Verletzung ist beim heutigen Stand der Medizin doch kein
Problem mehr}) oder Realativierungen durch zeitliche und
soziale Vergleiche (z.\,B. \Speech{\enquote{Im Vergleich zu den
anderen hatte ich großes Glück}}).}
{}
{}
{}
{}

\Text{Positive Selbstinstruktion}
{ Sie stellt eine Möglichkeit
zur emotionalen Selbstregulation dar. Gedanken,
Vorstellungsbilder und gedankliche Selbstgespräche können
innere Zustände verstärken oder abschwächen, im negativen
wie im positiven. Sich selbst zu beruhigen (\SpeechTextMargin{\enquote{Keine
Panik, nur die Ruhe}}), Mut zuzusprechen (\enquote{Du schaffst
das}), Trost zu spenden (\SpeechTextMargin{\enquote{Jeder macht mal
Fehler}}) kann helfen \Speech{\enquote{impulsive Reaktionen zu
unterdrücken [...], die eigene Aufmerksamkeit und Konzentration auf eine Aufgabe zu richten und
Frustration und Misserfolg zu kontrollieren}}
\cite{BengelPsyNotfallmedRd2-PsyBelRdPers}.}
{}
{}
{}
{}


\A{Gedankenstopp -  Imagination: Damit ist gemeint, dass man
die eigene Phantasie nutzt und sich hilfreiche Dinge
vorstellt. Wer seine Rettungsdienstjacke bewusst anzieht und
zuknöpft, kann sich z.\,B. gedanklich ausmalen, dass die
Jacke – wie eine Rüstung oder eine Art \enquote{Panzerung} – auch
einen psychischen Schutz bewirkt. Ebenso denkbar ist die
Vorstellung eines unsichtbaren Helfers oder Beschützers. In
diesem Zusammenhang können durchaus auch religiöse
Vorstellungen hilfreich sein. Sich aufdrängende negative
Gedanken lassen sich mit der Technik des Gedankenstopps
unterbrechen. Auch dabei handelt es sich letztlich um ein
imaginatives Verfahren.}

\A{Wahrnehmung lenken:  Vor besonders belastenden Anblicken
sollten Helfer sich schützen, indem sie sich – so banal dies
auch klingen mag – nach Möglichkeit abwenden und eben nicht
hinsehen. In bestimmten Situationen, etwa wenn sich jemand
in Anwesenheit von Einsatzkräften suizidiert, kann es sogar
angebracht sein, sich die Augen, ggf. auch die Ohren
zuzuhalten. Einer Person, die in Selbsttötungsabsicht vom
Dach eines Gebäudes springt, sollte man nicht hinterher
schauen usw. Wenn Situationen insgesamt als etwas sehr
Unangenehmes erlebt werden, sollte man sich auch nicht
länger am Ort des Geschehens aufhalten, als dies aus
medizinisch-einsatztaktischen Gründen unbedingt notwendig
ist. Extrem grausamen Anblicken (z.\,B. dem Anblick eines
stark entstellten Mordopfers oder eines bis zur
Unkenntlichkeit verbrannten Kindes) und ekelhaften Gerüchen
(z.\,B. Verwesungsgeruch) muss man sich, sobald der Tod des
Patienten zweifelsfrei festgestellt worden ist, ebenfalls
nicht länger aussetzen. Darüber hinaus wird empfohlen, nicht
in die Gesichter bzw. Augen von Verstorbenen zu schauen
(2).}

\A{Atmung, Entspannungstechniken:  Es gibt eine ganze Reihe
von Entspannungstechniken, die durchaus auch in einem
Einsatzgeschehen angewendet werden können und schon nach
kurzer Zeit wirksam sind. Eine sehr einfache Methode ist das
mehrfache tiefe Ein- und Ausatmen. Darüber hinaus können
z.\,B. Elemente bzw. einzelne Übungen der progressiven
Muskelrelaxation oder des Autogenen Trainings durchgeführt
werden. Diese Entspannungstechniken müssen zuvor jedoch in
entsprechenden Kursen erlernt und trainiert worden sein (7,
8).}

\EE{Ressourcen vor Augen halten}: Im Sinne des
Kohärenzgefühls bzw. der Theorie der Ressourcenerhaltung
kann es helfen sich die persönlichen, materiellen oder
sozialen Ressourcen bewusst zu machen um Anspannung zu
vermindern. \Speech{\enquote{So ist man bei einem Einsatz im
Rettungsdienst üblicherweise niemals allein. Die Kameraden sind ebenfalls
vor Ort, d.\,h. man arbeitet immer in einem Team. Zu den
materiellen Ressourcen gehört die persönliche
Schutzausrüstung, aber auch die Ausstattung des
Einsatzfahrzeuges. Im günstigsten Fall kann man sich sagen:
\Speech{´Ich bin gut gerüstet, jetzt kann es losgehen!´}}}
(Karutz, 2009).

\EE{Zeitliche Grenzen verdeutlichen}: Einsätze haben einen
zeitlichen Rahmen, d.\,h. sie haben einen Beginn, aber ebene
auch bei komplexeren Einsätzen (selbst Großschadenslagen) –
ein Ende. \Speech{\enquote{In der Regel sind Einsätze [\ldots]
innerhalb weniger Stunden abgeschlossen. Man ist den jeweiligen
Belastungen also nicht ewig, sondern lediglich für einen
definierten Zeitraum ausgesetzt. Dies sollte man sich
bewusst vor Augen führen}} \cite{Karutz200912}.

\EE{Zusätzliche Stressoren meiden}: Wenn die Belastung durch
das Notfalls geschehen ohnehin schon spürbar hoch ist, gilt
es zusätzliche Belastungsfaktoren zu minimieren bzw.
auszuweichen um das Erregungsniveau auf angenehmen bzw.
leistungsfähigem Level zu halten. \Study{Karutz}
\cite{Karutz200912} prägt hier den Begriff der \enquote{hilfreiche Langsamkeit}, d.\,h. \enquote{gerade wenn
man es eilig hat, eben nicht wie bei einem sportlichen
Wettkampf zu rennen}, zum Einsatzort zu \enquote{rasen}, den Funk im
Fahrzeug nicht unnötig laut einzustellen sowie Äußerungen zu
vermeiden, mit denen man Hektik erzeugt (z.\,B.
\SpeechTextMargin{\enquote{Los! Schneller!}}, \Speech{\enquote{Gas, Gas, Gas! Wir
müssen vor dem NEF da sein!}}.

\subsubsection{Langfristige Strategien}

\Text{Erholung}
{Sie dient dazu -- im Sinne der oben
dargestellten Stresskonzepte -- das homöostatische
Gleichgewicht wiederherzustellen, Widerstandsenergien erneut
aufzubauen, Ressourcen zu nutzen und zu erweitern,
symbolisch gesprochen also \enquote{die persönlichen Energietanks
wieder aufzufüllen} und auf neuerliche Stresssituationen
vorzubereiten. Während der Relaxation werden
Entspannungsreaktionen eingeleitet, die komplementär zur
Stressreaktion von physiologischen Prozessen begleitet
werden, nämlich dem Sinken der Muskelspannung, der
kortikalen Aktivität, der Herzfrequenz, des Blutdrucks und
der Atmung \cite{ZimbardoPsych6}. Auf
psychischer Ebene werden Wohlbefinden, Zufriedenheitsgefühle, Gefühle der
Gelassenheit sowie Selbstwahrnehmung gefördert. Unterstützt
kann dies durch div. Entspannungstechniken (progressiven
Muskelrelaxation, Autogenes Training, Mediatation,
Achtsamkeitsübungen etc.) werden, die sich tw. sogar während
des Einsatzgeschehens anwenden lassen (z.\,B.
Atemtechniken). Erholungsphasen sind von großer Bedeutung
und sollten dementsprechend bewusst gestaltet, in den Alltag
integriert und regelmäßig in angemessenem Ausmaß stattfinden
um ihre gesundheitserhaltende und -- fördernde Funktion
erfüllen zu können.}
{}
{}
{}
{}

\Text{Selbstwahrnehmung und Ausdruck von Emotionen}
{Wie
bereits zuvor erwähnt läuft die Entstehung von Stress nicht
immer bewusst, sondern sukzessive und unterschwellig ab.
Während des Einsatzes führen diverse Copingstrategien dazu,
Anspannung in Grenzen zu halten bzw. damit einhergehende
unangenehme emotionale Zustände (Ärger, Wut, Trauer etc.)
hintanzustellen bzw. zu unterdrücken, um handlungsfähig zu
bleiben und \enquote{kühlen Kopf zu wahren}. Diese zunächst
angestauten Affekte können sich aber zu einem weit späteren
Zeitpunkt im Stützpunkt oder daheim entladen, z.\,B. durch
Frustrationsgefühle bis hin zu aggressiven Verhalten
gegenüber anderen, ohne dass diese überschießenden
Reaktionen mit den ursprünglichen Ereignissen in Bezug
gesetzt werden. 

Die unmittelbarste Möglichkeit affektive
Erregungszustände abzuführen ist, sie durch gezielte
Selbstbeobachtung (in sich hören, Körperwahrnehmung) zu
identifizieren, innerlich zuzulassen und im Idealfall auch
auszudrücken. 
Bengel \cite{BengelPsyNotfallmedRd2-PsyBelRdPers} betont, dass gereizte-aggressive
unkontrollierte Emotionsäußerungen ineffektiv sind und es
wichtig wäre \enquote{unangenehme Gefühle in einer kontrollierten
und für die Umwelt sozial verträglichen Art und Weise
abzureagieren}. Eine Möglichkeit bestünde in Gesprächen mit
einem guten Freund, Angehörigen, Kollegen. Unspezifische
Frustrationen im Sinne von Anspannung, innerer Unruhe,
negativen Gedanken können auch durch das Ausüben
anstrengender Tätigkeiten, z.\,B. Sport, körperliche
Betätigung etc. abgebaut werden.}
{}
{}
{}
{}

\Text{Rituale}
{Sowohl am Arbeitsplatz als auch in Erholungsphasen können
ritualisierte Handlungen helfen, Entspannung zu fördern und
in den Alltag einzubinden. Das Tragen von Dienstkleidung als
bereits implementiertes Ritual soll helfen die Rollen als
professioneller Helfer und Privatperson zusätzlich zu
differenzieren, dies kann bewusst gemacht und verstärkt
werden durch das Hinlenken der Aufmerksamkeit während des
Umziehens -- \SpeechTextMargin{\enquote{Mit meiner Kleidung streife
ich auch alle Sorgen des heutigen Einsatztages ab -- jetzt habe ich
Freizeit}}. 

Gemeinsame Rituale wie z.\,B. morgendliche Teamsitzung mit
Frühstück, Einsatznachbesprechungen, gemeinsame
Kaffeepausen, etc. können nicht nur psychohygienischen
sondern auch teamfördernden Charakter haben und ermöglichen
einen informellen Austausch, auch im Sinne einer Möglichkeit
Gefühle, Sorgen, Ängste auszudrücken und Unterstützung zu
suchen. Rituale im Privaten z.\,B. ein Spaziergang nach der
Arbeit, der abendlich \Speech{\enquote{Wohlfühl-Tee}} oder der
Besuch des Lieblingsgeschäfts signalisieren evtl. sogar den Beginn von
Entspannungsphasen und können als konditionierter Reiz wird.}
{}
{}
{}
{}

\Text{Unterstützung durch die Umwelt}
{Zimbardo \cite{ZimbardoPsych6} verweist auf die Bedeutung
sozialer Unterstützung bei der Stressbewältigung und meint damit
\Speech{\enquote{Ressourcen, die von anderen Personen
bereitgestellt werden [...]}}. Beispielhaft zählt er
materielle Hilfen, emotionale Unterstützung (Fürsorge, Wertschätzung,
Sympathie, Zugehörigkeitsgefühl) und Zuwendung durch
Information (Ratschläge, Feedback) auf. Soziale Netze zu
entwickeln, zwischenmenschliche Beziehungen zu pflegen und
angebotene Hilfestellungen anzunehmen kann somit auf allen
Ebenen der Bewältigung wirksam sein: kurzfristig als
effektive Copingstrategie (konkrete Hilfe zur Bewältigung
einer Situation), als wertgeschätzte Ressource, zur
Förderung des Kohärenzgefühls (erweitertes
Handlungsrepertoire durch die Gemeinschaft) und im Rahmen
gemeinsamer Gestaltungsmöglichkeiten der Freizeit- und
Erholungsphasen.}
{}
{}
{}
{}


\A{Rationalisieren Zu dieser Strategie gehört, dass man die
Notfallsituation bzw. das Einsatzgeschehen möglichst
sachlich und nüchtern betrachtet. Medizinisch-technische,
insgesamt eher abstrakte Überlegungen rücken dabei in den
Vordergrund, während emotionale Aspekte bewusst
ausgeklammert werden. Zwar sind Notfallpatienten eben keine
defekten Maschinen, die mit \enquote{Vitalfunktionsmechanik}
lediglich \enquote{repariert} werden müssen, sondern Menschen, deren
spezifische psychische Situation bei der Hilfeleistung
angemessen zu berücksichtigen ist (6). In bestimmten
Einzelfällen (!) – nämlich dann, wenn ein Helfer ansonsten
handlungsunfähig würde – kann es allerdings sehr wohl
angebracht sein, z.\,B. ausschließlich auf \enquote{das Polytrauma},
\enquote{den Herzinfarkt} oder \enquote{die Reanimation} zu fokussieren.}

\A{Humor Gerade der Umgang mit  schwierigen Situationen kann
durch eine humorvolle Betrachtung erleichtert werden
(\enquote{Humor ist, wenn man trotzdem lacht}). Mitunter wird
Humor sogar explizit als Mechanismus zur Abwehr von
Stressreaktionen bezeichnet (2, 9). In diesem Zusammenhang
ist jedoch niemals Schadenfreude, Ironie oder Sarkasmus
gemeint, sondern Humor als Ausdruck der Fähigkeit, auch
schwierigen und bedrohlichen Situationen mit innerer
Distanz, Ruhe und Gelassenheit zu begegnen. In diesem Sinne
kann Humor sogar als Ausdruck von Hoffnung und Trost
verstanden werden. Banale Witzeleien erfüllen diese Funktion
eher nicht. Auch dass die humorvolle Betrachtung einer
Notfallsituation nicht mit Verhaltensweisen der Helfer
verbunden sein darf, die Notfallpatienten, Angehörige oder
andere Betroffene zusätzlich belasten, versteht sich
sicherlich von selbst.}

\A{Ablenkung Indem man sich auf die Durchführung einer
einzelnen Maßnahme konzentriert (z.\,B. auf das Vorbereiten
einer Infusion oder das Blutdruckmessen), können belastende
Aspekte des Notfallgeschehens zumindest vorübergehend
ausgeblendet werden. Auch leises Summen, das Aufsagen von
Schüttelreimen (\enquote{Fischers Fritz fischt frische Fische,
Frische Fische fischt Fischers Fritz}) oder das Durchgehen
des kleinen \enquote{Einmaleins} kann von belastenden Eindrücken
ablenken. Das bewusste Zählen oder Benennen von Gegenständen
im Umfeld hat ebenfalls eine solche Wirkung. Wenn es dazu
beiträgt, insgesamt handlungsfähig zu bleiben, kann sogar
ein kurzzeitiges \enquote{mentales Fremdgehen} legitim sein, d.\,h.
ein gedankliches \enquote{Aussteigen}, bei dem man für einen
Augenblick absichtlich an etwas anderes denkt.}

\A{Abb. 4: Eine
Möglichkeit der
psychologischen
Selbsthilfe:
Konzentration auf
die Durchführung
einer einzelnen
Maßnahme,
z.\,B. auf das
Vorbereiten einer
Infusion}

\A{Abb. 3: Die
meisten Helfer
\enquote{funktionieren}
auch in stressigen
Situationen und
benötigen an der
Einsatzstelle keine
psychologische
Unterstützung}

\A{Abb. 5: Auch
Kollegen können
helfen – mit
kleinen Gesten,
einem kurzen
Blickkontakt oder
einer aufmunternden
Äußerung}

\A{Die Auseinandersetzung mit den hier beschriebenen
Selbsthilfestrategien ist übrigens schon allein deshalb
sinnvoll, weil sie verdeutlicht, dass man auch bei starker
Belastung durchaus etwas tun kann, um einem drohenden
Verlust   der persönlichen Handlungsfähigkeit entgegen zu
wirken. Erwähnt werden soll auch, dass allein das Wissen um
solche Mechanismen hilfreich ist. Wer über
Selbsthilfestrategien  informiert ist, wird sie womöglich
schon deshalb nicht mehr    benötigen. Dies ist mit dem
Verfassen eines \enquote{Spickzettels} in der Schule vergleichbar:
Wer solch eine Merkhilfe erst einmal verfasst hat, kommt
erstaunlicherweise meist ohne sie aus.}

\A{Supervision:
Exkurs: Supervision - ein Tabu?
Abhängig von der Organisations- und Teamkultur kann
Supervision auch als erstes Anzeichen einer Überforderung,
Schwäche etc. fehlinterpretiert bzw. assoziiert werden.
Tatsächlich drückt Supervision ein gesundheitsförderliches
Verhalten und bewusst verantwortliches Engagement für den
Beruf aus, ein fairer Handel auf sein eigenes Wohlbefinden
und seine Gesundheit zu achten und in Austausch dafür der
Organisation seine unbeeinträchtigte berufliche
Leistungsfähigkeit verfügbar machen.}









\A{Stressbewältigung
Es gibt verschiedene Möglichkeiten der Stressreduktion und Stressbewältigung:
• Stressmanagement: Förderung kognitiver bzw. emotionsbezogene Copingstrategien
• Körperliche Stressbewältigung, z.\,B. durch Aerobes Training
• Entspannungsverfahren
• Humor
• Soziale Stressbewältigung: soziale Unterstützung
• Psychische Stressbewältigung, personale Ressourcen
• Anwendung von Widerstandsressourcen vor dem Hintergrund eines Kohärenzgefühl als Lebenseinstellung (ANTONOVSKY), d.\,h.
beispielsweise Probleme als Herausforderungen ansehen
• Resilienz
• Günstige Kontrollüberzeugungen (Attributionsstile)
• Hohe Selbstwirksamkeitserwartung
• Biofeedback}

\A{Keywords: Präventive, kurzfristige Strategien (Psychologische Selbsthilfe), Langfristige Strategien: personenenbezogen: Erholungsphasen, Selbsbeobachtung, Ablenkung, Rituale (z.\,B. aus der Rolle aussteigen), (Entlastung durch zwischenmenschliche Ressourcen, soziale Unterstützung), (Selbstbild vs. Ansprüche: Erkennen und Umgang mit eigenen Schwächen, Zulassen von Schwächen), (Dysfunktionale Strategien - Verhaltenssüchte); organisatorisch: Umgangsregeln im Teamkultur (Kommunikations- und Umgangsregeln),  Kohärenzgefühl durch Kompetenz (Fortbildung), (Supervision), (CISM). }

\A{Psychologische Grundbegriffe}












% \cite{AntonovskySalutogenese}:
% \bibentry{AntonovskySalutogenese} \\ 
% % Antonovsky, A. (1997). Salutogenese: zur
% % Entmystifizierung der Gesundheit. Dt. erw. Hrsg. von Alexa Franke. Deutsche
% % Gesllschaft für Verhaltenstherapie. Tübingen: Dgvt-Verl.
% 
% \cite{BengelPsyNotfallmedRd2-PsyBelRdPers}:
% \bibentry{BengelPsyNotfallmedRd2-PsyBelRdPers} \\ 
% % Bengel, J. \& Heinrichs, M. (2004). Psychische
% % Belastungen des Rettungspersonals. In J. Bengel (Hrsg.), Psychologie in
% % Notfallmedizin und Rettungsdienst (2. Auflage; S. 25-43).
% % Berlin: Springer.
% 
% \cite{BergnerBurnOutAerzten9}:
% \bibentry{BergnerBurnOutAerzten9} \\ 
% % Bergner, Thomas (2004). Burn-out bei Ärzten:
% % Lebensaufgabe statt Lebens-Aufgabe, 9,  S. 410.
% 
% \cite{BurischBurnOut2}: \bibentry{BurischBurnOut2} \\
% % Burisch, M. (1994). Das Burnout-Syndrom – Theorie
% % der inneren Erschöpfung (2. Auflage). Berlin: Springer.
% 
% \cite{BurischBurnOut3}: \bibentry{BurischBurnOut3} \\
% % Burisch, M. (2006). Das Burnout Syndrom – Theorie
% % der inneren Erschöpfung (3., überarbeitete Auflage). Berlin:
% % Springer.
% 
% \cite{DavisonKlinPsy1988}: \bibentry{DavisonKlinPsy1988} \\
% % Davison, G. C. \& Neale, J. M. (1988): Klinische
% % Psychologie. 3. Auflage. München: PVU. Deutsches Institut
% % für Medizinische Dokumentation und Information (2011).
% 
% \cite{ICD10-De-2011}
% 
% \cite{Freudenberger1974}: \bibentry{Freudenberger1974} \\
% % Freudenberger, H. J. (1974). Staff Burn-Out. Journal of
% % Social Issues, Vol. 30 (1), 159–165.
% 
% \cite{FrielingLehrbArbeitspsy2}:
% \bibentry{FrielingLehrbArbeitspsy2} \\ 
% % Frieling E. \& Sonntag Kh. (1999). Lehrbuch
% % Arbeitspsychologie (2. vollst. überarb. Und erweiterte
% % Aufl.). Bern: Huber.
% 
% \cite{PlutchikEmotionVol3-NeuroendocrinePatterns}:
% \bibentry{PlutchikEmotionVol3-NeuroendocrinePatterns} \\
% % Henry, J. P. (1986). Neuroendocrine patterns of emotional response. In Plutchik H. \& Kellerman H. (Hrsg.), Emotion:
% % Theory, Research and Experiences (Vol. 3; S. 37-60), San
% % Diego: Academic Press.
% 
% \cite{Karutz200912}: \bibentry{Karutz200912} \\
% % Karutz, H. (2009). Wenn die Belastungsgrenze erreicht ist:
% % Psychologische Selbsthilfe in Extremsituationen.
% % Rettungsdienst, 12, 40-45.
% 
% \cite{NitschStress-StressTransakt}:
% \bibentry{NitschStress-StressTransakt} \\ 
% % Lazarus, R. S. \& Launier, R (1981). Stressbezogene
% % Transaktion zwischen Person und Umwelt. In Nitsch, J. R.
% % (Hrsg.), Stress, Theorien, Untersuchungen, Maßnahmen (S.
% % 213-260). Bern: Huber.
% 
% \cite{LazarusStress}: \bibentry{LazarusStress} \\
% % Lazarus, R.S. \& Folkman, S. (1984). Stress, appraisal, and
% % coping. New York: Springer.
% 
% \cite{MaslachBurnout}: \bibentry{MaslachBurnout} \\
% % Maslach, C. (1982). Burnout. The Cost of Caring. Englewood
% % Cliffs, NJ: Prentice-Hall.
% 
% \cite{OskampAppSocPsyAnn84-BurnOrgSett}:
% \bibentry{OskampAppSocPsyAnn84-BurnOrgSett} \\ 
% % Maslach, C. \& Jackson, S. E. (1984). Burnout in
% % organizational settings. In S. Oskamp (Hrsg.), Applied
% % Social Psychology Annual (S. 133-153). Beverly Hills, CA:
% % Sage.
% 
% \cite{Niaura2002}: \bibentry{Niaura2002} \\
% % Niaura, R., Todaro, J. F., Stroud, L., Spiro, A., Ward, K.
% % D., \& Weiss, S. (2002). Hostility, the metabolic syndrome,
% % and incident coronary heart disease. Health Psychology, Vol.
% % 21(6), 588-593. 
% 
% 
% \cite{StarkeMedSozAspProbBew00}: \bibentry{StarkeMedSozAspProbBew00} \\
% % Starke, D. (2000). Medizinsoziologie. Kognitive, emotionale
% % und soziale Aspekte menschlicher Problembewältigung: ein
% % Beitrag zur aktuellen Stressforschung. Münster: LIT.
% 
% \cite{BruederlBelastung88-DefBew}: \bibentry{BruederlBelastung88-DefBew} \\
% % Trautmann-Sponsel, R. D. (1988): Definition und Abgrenzung
% % des Begriffs Bewältigung. In L. Brüderl (Hrsg.), Belastung
% % und Bewältigung: Trends in der Bewältigungsforschung (S.
% % 14--24). München: Juventa.  
% 
% \cite{Yerkes1908}: \bibentry{Yerkes1908} \\
% % Yerkes, R.M. \& Dodson, J.D (1908). The relation of strength
% % of stimulus to rapidity of habit-formation. Journal of
% % Comparative Neurology and Psychology, 18, 459-482.
% 
% \cite{ZimbardoPsych6}: \bibentry{ZimbardoPsych6} \\
% % Zimbardo, G.  (1995). Psychologie (6. Auflage). Berlin:
% % Springer.

\cite{AntonovskySalutogenese,BengelPsyNotfallmedRd2-PsyBelRdPers,
BergnerBurnOutAerzten9, BurischBurnOut2, BurischBurnOut3,
DavisonKlinPsy1988, ICD10-De-2011, Freudenberger1974,
FrielingLehrbArbeitspsy2,
PlutchikEmotionVol3-NeuroendocrinePatterns, Karutz200912,
NitschStress-StressTransakt, LazarusStress, MaslachBurnout,
OskampAppSocPsyAnn84-BurnOrgSett, Niaura2002,
StarkeMedSozAspProbBew00, BruederlBelastung88-DefBew,
Yerkes1908, ZimbardoPsych6}



% Selye, H. (1956). The Stress of life. New
% York: McGrw-Hill. In P. G. Zimbardo, (Hrsg.), Psychologie
% (6. Auflage; S. 579-580). Berlin: Springer.







%%%%%%%%%%%%%%%%%%%%%%%%%%%%%%%%%%%%%%%%%%%%%%%%%%%%%%%%%%%
%%%%%%%%%%%%%%%%%%%%%%%%%%%%%%%%%%%%%%%%%%%%%%%%%%%%%%%%%%%
%%%%%%%%%%%%%%%%%%%%%%%%%%%%%%%%%%%%%%%%%%%%%%%%%%%%%%%%%%%
%%%%%%%%%%%%%%%%%%%%%%%%%%%%%%%%%%%%%%%%%%%%%%%%%%%%%%%%%%%
%%%%%%%%%%%%%%%%%%%%%%%%%%%%%%%%%%%%%%%%%%%%%%%%%%%%%%%%%%%
%%%%%%%%%%%%%%%%%%%%%%%%%%%%%%%%%%%%%%%%%%%%%%%%%%%%%%%%%%%
\section{Psychische Betreuung von Kranken/Verletzten}


\Text{{\TxQuerverweise}}{%
\begin{itemize}
  \item Gesprächsführung, Beziehungsaufbau: \dotfill
  \FullPrettyRef{topic:kommunikationsregeln}
  \item Begleitung Sterbender: \dotfill
  \FullPrettyRef{topic:sterbende-betreuung}
  \item Umgang mit psychiatrischen Patienten:
  \dotfill \FullPrettyRef{topic:umgang-psy}
\end{itemize}
}{}{}{}{}    
%     \Text{Beziehungsaufbau generell}
% {
% 
% }
% {}
% {}
% {}
% {}
% 
% 
% \Text{Suchterkrankte Menschen}
% {
% 
% }
% {}
% {}
% {}
% {}





% %%%%%%%%%%%%%%%%%%%%%%%%%%%%%%%%%%%%%%%%%%%%%%%%%%%%%%%%%%%
% %%%%%%%%%%%%%%%%%%%%%%%%%%%%%%%%%%%%%%%%%%%%%%%%%%%%%%%%%%%
% %%%%%%%%%%%%%%%%%%%%%%%%%%%%%%%%%%%%%%%%%%%%%%%%%%%%%%%%%%%
% %%%%%%%%%%%%%%%%%%%%%%%%%%%%%%%%%%%%%%%%%%%%%%%%%%%%%%%%%%%
% %%%%%%%%%%%%%%%%%%%%%%%%%%%%%%%%%%%%%%%%%%%%%%%%%%%%%%%%%%%
% \subsection{verwirrte Patienten}
% 
%     wird integrativ in allen Kapitel abgehandelt
%     




%%%%%%%%%%%%%%%%%%%%%%%%%%%%%%%%%%%%%%%%%%%%%%%%%%%%%%%%%%%
%%%%%%%%%%%%%%%%%%%%%%%%%%%%%%%%%%%%%%%%%%%%%%%%%%%%%%%%%%%
%%%%%%%%%%%%%%%%%%%%%%%%%%%%%%%%%%%%%%%%%%%%%%%%%%%%%%%%%%%
%%%%%%%%%%%%%%%%%%%%%%%%%%%%%%%%%%%%%%%%%%%%%%%%%%%%%%%%%%%
%%%%%%%%%%%%%%%%%%%%%%%%%%%%%%%%%%%%%%%%%%%%%%%%%%%%%%%%%%%







\Experimental{
%%%%%%%%%%%%%%%%%%%%%%%%%%%%%%%%%%%%%%%%%%%%%%%%%%%%%%%%%%%
%%%%%%%%%%%%%%%%%%%%%%%%%%%%%%%%%%%%%%%%%%%%%%%%%%%%%%%%%%%
%%%%%%%%%%%%%%%%%%%%%%%%%%%%%%%%%%%%%%%%%%%%%%%%%%%%%%%%%%%
%%%%%%%%%%%%%%%%%%%%%%%%%%%%%%%%%%%%%%%%%%%%%%%%%%%%%%%%%%%
%%%%%%%%%%%%%%%%%%%%%%%%%%%%%%%%%%%%%%%%%%%%%%%%%%%%%%%%%%%
%%%%%%%%%%%%%%%%%%%%%%%%%%%%%%%%%%%%%%%%%%%%%%%%%%%%%%%%%%%
\section{Supervision}

    wird extern abgehandelt (PEER-Vortrag)









}













